\documentclass[12pt]{report}

% --------------------
% Basic packages
% --------------------
\usepackage[T1]{fontenc}
\usepackage{lmodern}
\usepackage{setspace}
\usepackage{geometry}
\usepackage{microtype}
\usepackage[hidelinks]{hyperref}
\usepackage{parskip}
\usepackage{amsmath}
\usepackage{csquotes}
\usepackage{enumitem}

% -------------------
% Draft Watermark
% -------------------

\usepackage{xcolor}

% --------------------
% Bibliography
% --------------------

\usepackage[
  backend=biber,
  style=authoryear,
  citestyle=authoryear
]{biblatex}

\addbibresource{thesis_bibliography.bib}

% --------------------
% Page layout
% --------------------
\geometry{
  margin=1in
}
\onehalfspacing
\setlength{\parskip}{1\baselineskip}

% --------------------
% Metadata
% --------------------
\title{The Invisible Limits of Public Reason:\\
Stability, Tracking, and the Political Epistemology of Justification}
\author{Alex Horovitz}
\date{\today}

% --------------------
% Document
% --------------------
\begin{document}
\pagenumbering{roman} 

\maketitle

\tableofcontents
\clearpage
\pagenumbering{arabic}

% ====================
% Introduction
% ====================

\chapter{Introduction: Legitimacy under Drift}

\section{The puzzle: regimes of justification in a changing environment}

In 1989, it was common to hear that ``western history'' had reached its endpoint in the ``triumph'' of liberal democracy.  Yet, within a generation, liberal democratic institutions became ``fragile in most people's living memory,'' and their stability again appeared as a central political problem \parencite[p.1187]{Weithman2024Rescuing}.  The worry is not only that democracies may fail to deliver desired outcomes, but that they may fail in a more basic way: as regimes that claim \emph{legitimacy} through public justification rather than mere command.

Rawlsian political liberalism is the canonical attempt to explain how legitimacy is possible under conditions of deep and enduring moral diversity.  Rawls frames one of political philosophy's most important aims, in a constitutional democracy, as presenting a political conception of justice that can (i) provide ``a shared public basis for the justification of political and social institutions'' and (ii) ``help ensure stability from one generation to the next''.  A justification that rests only on ``self- or group-interests'' cannot do that work; even if institutional design moderates conflict, it yields at best a ``mere \emph{modus vivendi}'' dependent on contingency \parencite[p.~1]{rawls1987overlapping}.  The aspiration, then, is stability \emph{for the right reasons}: stability grounded in citizens' endorsement of a publicly justifiable political conception, not in temporary alignment of forces.

The puzzle that motivates this thesis is that contemporary politics rarely provides a stable backdrop for that aspiration.  Even when citizens are reasonable and sincerely disposed to cooperate, democratic politics takes place in environments that change: problems evolve, technologies reshape the costs of policies, and institutions revise the meaning and application of constitutional principles.  Complexity theory offers a useful image for this condition.  In many domains, the ``payoff'' of a policy depends on what others are doing, so the evaluative terrain itself shifts as agents adapt.  Page calls these coupled, co-evolving terrains ``\emph{dancing landscapes}'' \parencite[p.~94]{page2011diversity}.  If the landscape dances, then a justification that succeeds at time $t$ may fail at time $t+\Delta t$, not because citizens became unreasonable, but because the mapping from shared political values to concrete consequences changed.

Rawlsian public reason is designed to answer a question that this dynamism makes acute: what kinds of reasons may citizens reasonably offer one another when \emph{fundamental political questions} are at stake?  Rawls begins \emph{The Idea of Public Reason Revisited} by emphasizing that public reason ``belongs to a conception of a well ordered constitutional democratic society'' and that its form and content are ``part of the idea of democracy itself'' \parencite[p.~765]{rawls1997public}.  The central pressure arises from the fact of reasonable pluralism: a plurality of incompatible yet reasonable comprehensive doctrines is the ``normal result'' of free institutions.  Citizens therefore cannot expect to converge, even in principle, on the basis of their comprehensive views; they must instead consider what reasons they can offer one another as \emph{citizens} \parencite[p.~766]{rawls1997public}.  The ambition of public reason is to secure legitimacy \emph{in spite of} this persistent pluralism.

The central question of this thesis is whether that ambition is \emph{feasible} in the political environments we inhabit: environments marked by large-scale mediation, high informational complexity, and rapid institutional and technological change.

\section{Contribution and target}

This dissertation targets a specific legitimating route within Rawlsian political liberalism: \textbf{legitimacy through stability-by-public-justification}.  The focus is not Rawls's moral values in abstraction, nor a general dismissal of deliberation or democratic contestation.  It is a test of whether Rawls's stability strategy can function as advertised when the social and epistemic conditions needed for it are made explicit.

Two features of Rawls's own presentation sharpen the target.

First, Rawls insists that the public conception of justice in a constitutional democracy should be ``so far as possible, independent of controversial philosophical and religious doctrines,'' applying ``the principle of toleration to philosophy itself'' \parencite[p.~223]{rawls1985justice}.  Political liberalism seeks a justificatory basis that citizens with diverse comprehensive doctrines can nevertheless endorse.

Second, Rawls emphasizes that political liberalism is a \emph{freestanding} political doctrine: it ``works entirely within'' the political domain and ``does not rely on anything outside it.'' Political conceptions of justice are presented as political, not as consequences of any particular comprehensive doctrine \parencite[p.~133]{rawls1995reply}.  The aspiration is legitimacy through reasons that are public---reasons that can be shared by free and equal citizens in their political relationship.

This thesis argues that this legitimating route has \emph{invisible feasibility conditions}.  Very roughly: for stability-through-public-reason to be more than episodic, citizens must be able to \emph{track} justificatory status across time, and the justificatory ``landscape'' must be stable enough for that tracking to succeed.  These conditions are rarely stated explicitly in Rawlsian theory, yet they do crucial work.  Once we foreground them, a distinctive vulnerability appears: even if citizens are reasonable in Rawls's normative sense, they may be unable to perform the diachronic epistemic tasks that stability-through-reasons presupposes \emph{under drift}.

To avoid a common misunderstanding, it is important to keep Rawls's normative concept of \emph{the reasonable} separate from a psychological claim about what agents can do.  Rawls emphasizes that reasonable pluralism is not a defect of citizens' character but a normal feature of democratic culture \parencite[p.~766]{rawls1997public}.  The feasibility question raised here is different: given that pluralism is persistent, and given the complexity and dynamism of modern political environments, can citizens \emph{reliably} satisfy the epistemic burdens implied by stability-through-public-justification?

\section{Two conclusions, kept separate}

The argument of this dissertation yields two conclusions that must be kept distinct.

\subsection*{(A) The feasibility claim}

Even if Rawlsian political liberalism identifies genuine moral values, there may be \emph{no feasible institutional form} that delivers stability-through-shared-public-justification at scale under real conditions of drift, mediation, and complexity.  The point is not that democratic justification is always impossible, but that the specific Rawlsian route---legitimacy through citizens' stable endorsement of publicly justified coercive law---depends on capacities and environmental regularities that may not obtain.

\subsection*{(B) The internal-critique claim}

If Rawlsian legitimacy is defined by the stability-through-reasons ideal, then there are \emph{limits internal to the ideal} that become visible once we foreground its tracking and landscape assumptions.  Feasibility constraints here are not merely external ``realist'' complaints; they may delimit the domain in which this particular account of legitimacy can do the work it is assigned.  Rawls himself treats stability as part of political philosophy's aim in constitutional democracy \parencite[p.~1]{rawls1987overlapping}, and recent work underscores that stability remains politically urgent in ``our current political moment'' \parencite[p.~1187]{Weithman2024Rescuing}.  If stability is doing legitimating work, then the conditions under which it is achievable are not a side issue: they shape what legitimacy via public reason can plausibly be.

The upshot is not anti-democratic skepticism.  Nor is it a claim that ``citizens are irrational'' or that public deliberation is worthless.  Rather, it is a diagnosis of a potential mismatch between (i) what stability-through-reasons demands and (ii) what bounded agents in drifting environments can do.  If that mismatch is real, then either the Rawlsian story about the citizen's role must be revised, or institutions must shoulder more of the epistemic labor while preserving contestability and civic respect.

\section{Timescale preview and a vignette}

A key theme of the dissertation is that feasibility depends on \emph{timescale}.  ``Stability'' can be demanded over different horizons, and citizens' justificatory updating has different prospects at different speeds of change.  Two rough regimes matter throughout:

\begin{itemize}
  \item \textbf{Short-cycle drift (months to a few years):} agenda shocks, crises, rapid technological change, and institutionally significant reinterpretations.
  \item \textbf{Long-cycle drift (decade+):} generational replacement, slow evolution of doctrines and identities, and path-dependent institutional change.
\end{itemize}

Rawls's idea of public reason is explicitly keyed to fundamental political questions: its subject is the ``public good concerning questions of fundamental political justice,'' in particular ``constitutional essentials and matters of basic justice'' \parencite[p.~767]{rawls1997public}.  A tempting response to feasibility worries is therefore a \emph{granularity} reply: perhaps citizens need not track fine-grained policy reasons; perhaps it is enough to track the justification of constitutional essentials in a coarse way.  This dissertation takes that reply seriously (especially in its strongest contemporary forms).  But timescale and drift matter even at coarse levels, because constitutional concepts themselves are interpreted through evolving institutions, technologies, and social facts.

A brief vignette illustrates the tracking problem without trying to prove it empirically.  Consider a familiar triad: privacy, speech, and security.  A state may justify expanded surveillance powers by appeal to public goods like security and the protection of basic liberties.  Yet as surveillance technologies improve, as threats change, and as institutional safeguards erode or strengthen, the same legal category (``metadata,'' ``reasonable suspicion,'' ``public forum'') can come to have very different consequences for citizens' liberties.  The justificatory landscape thus shifts with the co-evolution of policy, technology, and behavior; a dynamic well captured by Page's notion of ``dancing landscapes.''  If the landscape shifts quickly, citizens face a recurring task: to determine whether the coercive scheme remains justifiable to them and to others on public grounds, and to update their political support without treating every change as a wholesale reset.

The later chapters formalize this intuitive picture.  They argue that stability-through-reasons implicitly requires both (i) a citizen-facing capacity to track justificatory status across time and (ii) sufficient stability in the mapping from reasons and facts to coercive outcomes.  The argument then asks when those conditions are likely to fail, and what legitimacy should look like if they do.

\section{Roadmap}

Chapter~2 reconstructs the minimal Rawlsian architecture needed for the critique: stability for the right reasons, public reason, and the distinction between the \emph{reasonable} and the \emph{rational}.  Chapter~3 situates the project in the feasibility literature and explains why feasibility is an \emph{internal} constraint for a legitimacy route that relies on stability-by-justification.  Chapter~4 defines the two ``invisible'' assumptions---tracking and landscape stability---and introduces the timescale and granularity parameters that organize the later debate.

Chapters~5--7 develop the two main sources of failure: (i) citizen-side epistemic burdens under mediation and bounded cognition, and (ii) world-side nonstationarity in complex adaptive political environments.  Chapter~8 gives the strongest Rawlsian recovery strategies, including coarse-graining and set-valued stability replies, and argues that these strategies only partially resolve the problem.  Chapter~9 turns constructive: if individual high-resolution tracking is too demanding, legitimacy may need publicly inspectable institutional scaffolds that reduce cognitive load while preserving contestability.  Chapter~10 concludes by drawing implications for political liberalism, deliberative democracy, and institutional design under persistent pluralism and drift.

\chapter{Rawlsian Stability and Public Reason}

Chapter~1 framed the thesis around a practical problem for liberal legitimacy: modern political life is volatile, mediated, and interpretively plastic, yet liberal democracies often present themselves as regimes of justification rather than mere regimes of command. This chapter supplies the minimum Rawlsian machinery needed to diagnose that problem as an \emph{internal} challenge to political liberalism. The aim is not a full reconstruction of Rawls’s system, but a disciplined account of (i) \emph{stability for the right reasons}, (ii) the structure of \emph{public reason}, and (iii) the motivational picture that requires citizens to be both \emph{reasonable} and \emph{rational}. I close by introducing Paul Weithman’s influential interpretation of Rawlsian stability as a form of equilibrium, and by explaining why Rawls’s own late “pluralism about justice” concession pushes that equilibrium away from a single fixed point and toward a set-valued conception.

\section{Stability for the Right Reasons: What’s at Stake}

Rawls’s mature political project is organized around a demanding ambition: a constitutional democracy should be able to offer a \emph{shared public basis} for justifying its coercive institutions, and that basis should sustain the regime \emph{over time}, “from one generation to the next.” The ambition is not simply to identify principles that would be just \emph{if} complied with, but to describe a public political conception that can function as the public charter of a democratic people.

Two contrasts make Rawls's concern with stability more precise. First, he distinguishes stable social unity from a merely strategic compromise. A political arrangement resting on ``self- or group-interests alone'' cannot reliably reproduce itself; even if it is constitutionally well-designed, it remains a modus vivendi dependent on shifting contingencies.\parencite[p.~1]{rawls1987overlapping}

Rawls later makes the same point by emphasizing that stability depends on the moral psychology and doctrinal resources of citizens. Whether a political conception will be stable “importantly depends on the content of the religious, philosophical, and moral doctrines available to constitute an overlapping consensus.”\parencite[p.~250]{rawls1985justice} He illustrates the thought by imagining citizens who endorse a political conception (such as justice as fairness) from within very different comprehensive outlooks: a religious view that underwrites toleration, a comprehensive liberal moral view (Kantian or Millian), or a view that treats the political conception as itself sufficient in a democratic context.\parencite[pp.~250--251]{rawls1985justice} The stability generated by these diverse routes of endorsement is, Rawls insists, different in kind from the stability of a modus vivendi: it is not simply that citizens are \emph{willing} to comply so long as others do, but that they can see the political conception as expressing political values that normally outweigh competing considerations within their own outlooks.\parencite[p.~250]{rawls1985justice}

This stability has a distinctive normative shape in political liberalism. In his reply to Habermas, Rawls describes the “basis of social unity” in a well-ordered constitutional democracy as a condition in which (a) the basic structure is regulated by “the most reasonable political conception of justice,” (b) that conception is “endorsed by an overlapping consensus” of reasonable comprehensive doctrines, and (c) public political discussion about constitutional essentials and basic justice is “reasonably decidable” by the reasons supplied by that conception (or by “a reasonable family” of such conceptions).\parencite[p.~147]{rawls1995reply} When these conditions are satisfied, Rawls claims, citizens’ deepest convictions can support the political conception, and “\emph{from this follows stability for the right reasons}.”\parencite[p.~147]{rawls1995reply} When they are not satisfied, the regime’s unity is merely strategic: “there is only a modus vivendi,” and stability depends on a shifting balance of forces.\parencite[p.~147]{rawls1995reply}

The thesis will argue that this picture hides feasibility conditions---conditions about what citizens can reliably do, and about how stable the justificatory environment must be for that doing to succeed. But the critical point for now is structural: for Rawls, stability is not a cosmetic sociological add-on. Stability is part of the point of the political conception, and overlapping consensus is the mechanism that is meant to secure it \emph{without} converting legitimacy into mere bargaining or domination.\parencite[p.~1]{rawls1987overlapping}

\section{The Three Parts of Public Reason}

Rawls introduces public reason as a response to a basic democratic predicament: citizens must decide fundamental political questions despite deep, permanent disagreement about comprehensive moral and religious truth. In \emph{The Idea of Public Reason Revisited}, he says the idea “belongs to a conception of a well ordered constitutional democratic society,” and that its “form and content” are “part of the idea of democracy itself.”\parencite[p.~765]{rawls1997public} Public reason is meant to specify the kinds of reasons citizens may offer one another when the exercise of political power must be justified to all who are subject to it.

Rawls also insists that public reason has a \emph{definite structure}. He lists five aspects: the political questions to which it applies; the persons to whom it applies (government officials and candidates for public office); the content of public reason (given by a family of reasonable political conceptions of justice); the application of those conceptions to coercive law; and citizens’ checking that those principles satisfy a “criterion of reciprocity.”\parencite[p.~767]{rawls1997public} For the purposes of this thesis, it is useful to compress these into three parts:

\begin{enumerate}
  \item \textbf{Content:} a restricted body of political values and principles that can be expressed by a “family of reasonable conceptions of political justice,” each of which is reasonably thought to satisfy reciprocity.\parencite[p.~767]{rawls1997public}
  \item \textbf{Subject:} the fundamental questions to which public reason applies, which Rawls divides into “constitutional essentials” and “matters of basic justice.”\parencite[p.~767]{rawls1997public}
  \item \textbf{Forum:} the distinctive public arena in which those reasons must be used: the “public political forum,” divided into judicial discourse (especially a supreme court), the discourse of government officials (executives and legislators), and the discourse of candidates and their campaign representatives.\parencite[p.~767]{rawls1997public}
\end{enumerate}

These three parts matter because they locate Rawls’s justificatory demands. Public reason is not meant to govern every mode of democratic speech. Rawls sharply distinguishes the public political forum from the “background culture,” the diverse associational and intellectual life of civil society, and he explicitly denies that public reason applies to that background culture “nor to media of any kind.”\parencite[p.~768]{rawls1997public} In the background culture, “full and open discussion” is appropriate; political liberalism “fully agrees.”\parencite[p.~768]{rawls1997public}

At the same time, Rawls does not treat public reason as a purely professional ethic for officials. In addition to the officials and candidates who must directly use public reason in the forum, Rawls imagines ordinary citizens as playing a supervisory and participatory role: when constitutional essentials or basic justice are at stake, citizens should ideally regard themselves “as if they were legislators following public reason,” and (under that idealization) majority enactments can be legitimate even when they are not each citizen’s most favored law.\parencite[p.~771]{rawls1997public} The point is that legitimacy depends on the \emph{mode} of justification, not on unanimity.

\section{Burdens of Judgment and the Persistence of Reasonable Pluralism}

Public reason is motivated by Rawls’s diagnosis that deep disagreement is not a temporary defect to be cured by education, modernization, or improved argumentation. A “plurality of conflicting reasonable comprehensive doctrines,” he says, is “the normal result” of a democratic culture “of free institutions.”\parencite[p.~766]{rawls1997public} Citizens therefore cannot expect to resolve fundamental political questions by appealing to their own comprehensive doctrines and demanding that others submit. Instead, they “need to consider what kinds of reasons they may reasonably give one another when fundamental political questions are at stake.”\parencite[p.~766]{rawls1997public}

The engine of this diagnosis is Rawls’s idea of the \emph{burdens of judgment}: the predictable sources of reasonable disagreement among people who are sincere, competent, and willing to reason in good faith. While Rawls’s fullest discussion is in \emph{Political Liberalism}, he summarizes the core idea in \emph{Public Reason Revisited}: conflicts “arising from the burdens of judgment always exist and limit the extent of possible agreement,” and reasonable persons are partly characterized by their willingness to “recognize and accept the consequences of the burdens of judgment.”\parencite[p.~805]{rawls1997public} In his reply to Habermas, Rawls builds the same recognition directly into the concept of the reasonable: to be reasonable is, in part, to recognize these burdens and to accept their consequences for one’s attitude of toleration toward other comprehensive doctrines.\parencite[p.~134]{rawls1995reply}

Two clarifications are crucial for what follows. First, Rawls treats pluralism as a \emph{feature} of free democratic life rather than a pathology. Public reason is designed to respect the fact that citizens’ deepest commitments may be irreconcilable, while still making civic justification possible.\parencite[p.~766]{rawls1997public} Second, the fact of pluralism is not merely synchronic: it is not only that many doctrines coexist at once, but also that democratic societies reproduce themselves through changing generations, institutions, and interpretive traditions. Rawls’s original motivation for overlapping consensus is explicitly diachronic---it is meant to secure stability “from one generation to the next.”\parencite[p.~1]{rawls1987overlapping} Later chapters will use this diachronic dimension to frame the core feasibility question: whether citizens can keep justificatory assessments stable as doctrines, institutions, and political circumstances evolve.

\section{Accessibility and the “Public” Condition}

The public character of public reason is not exhausted by the sociological fact that citizens happen to share some values. It is instead a normative constraint on the \emph{reasons} offered in justification of coercive law. Rawls’s key device for articulating this constraint is the \emph{criterion of reciprocity}. Citizens are reasonable, Rawls says, when they are prepared to propose fair terms of cooperation and to act on them, even when this is costly, provided others do likewise; and the criterion of reciprocity adds a further condition on proposals of fundamental terms: those who propose the terms “must also think it at least reasonable for others to accept them, as free and equal citizens, and not as dominated or manipulated, or under the pressure of an inferior political or social position.”\parencite[p.~771]{rawls1997public}

This reciprocity constraint supports Rawls’s conception of political legitimacy. He writes that legitimacy “based on the criterion of reciprocity” holds that the exercise of political power is proper only when we sincerely believe that the reasons we would offer in our public role are sufficient, and when we “reasonably think that other citizens might also reasonably accept those reasons.”\parencite[p.~771]{rawls1997public} In other words, legitimate law must be justifiable in terms that others can acknowledge as free and equal co-authors of the political order, even when they do not share one another’s comprehensive views.

It matters, again, that this is not a demand for universal agreement. Rawls explicitly allows that citizens will differ about which political conceptions are most reasonable, and that laws enacted by majority procedures may fail to be each citizen’s preferred law.\parencite[p.~771]{rawls1997public} The demand is that the justificatory currency in the public political forum be political and shareable in principle---a demand that Rawls expresses by proposing that comprehensive doctrines of truth or right be “replaced by an idea of the politically reasonable addressed to citizens as citizens” when fundamental political questions are at stake.\parencite[p.~766]{rawls1997public}

Finally, Rawls’s restriction to the public political forum shows why the accessibility condition is simultaneously narrower and more demanding than it first appears. It is narrower because it does not govern the whole of civil society; it is more demanding because it targets precisely those moments---constitutional essentials and matters of basic justice---when the exercise of coercive power is most consequential and least avoidable.\parencite[p.~767]{rawls1997public} Later chapters will argue that satisfying accessibility over time requires citizens (and institutions) to maintain reliable, publicly intelligible markers of what counts as an admissible public justification.

\section{Rational and Reasonable}
The thesis's feasibility critique can easily be misread as assuming selfishness or bad faith: as if public reason fails because citizens are too irrational, too biased, or too unwilling to cooperate. To prevent that confusion, it is important to mark a distinction Rawls insists on: the difference between the \emph{rational} and the \emph{reasonable}.

Rawls treats both as independent aspects of practical reason. The rational concerns a person's good and the pursuit of a rational plan of life---Rawls labels his value theory ``goodness as rationality'' and develops it as a freestanding account of what it means to form and pursue ends well.\parencite[pp.~347,~359--365]{rawls1999theory} The reasonable, by contrast, is not derivable from rationality alone. Rawls identifies two core elements: first, the willingness to propose fair terms of social cooperation that others, as free and equal, might also endorse, and to honor those terms when others act likewise; second, recognition of the burdens of judgment and their consequences for toleration.\parencite[p.~134]{rawls1995reply} A perfectly rational person can still be unreasonable, acting in ways that violate independent moral constraints; conversely, a person can be reasonable even when reasonableness yields no private advantage.

Samuel Freeman's exposition illuminates why this matters for political theory: Rawls's constructivist picture presupposes citizens ``who are morally motivated by their sense of justice,'' yet whose agreement must also be compatible with the constraints of social cooperation and human motivation.\parencite[p.~xi]{freeman2007rawls} The political conception must speak to both dimensions at once.

Why does Rawlsian stability depend on both? The political conception is meant to regulate a society of free and equal citizens who pursue diverse ends and inevitably encounter conflicts of interest. Stability requires that citizens can see the political order as compatible with their good (a rational concern) while also being prepared to restrain their claims by fair terms of cooperation (a reasonable concern). In \emph{Public Reason Revisited}, Rawls describes reasonable citizens as those prepared to offer fair terms of cooperation and to act on them even when doing so runs against their interests in particular situations.\parencite[p.~771]{rawls1997public} The stability Rawls seeks is therefore not mere modus vivendi; it depends on citizens who are both rationally motivated and genuinely reasonable.

The point for this thesis is that even if we grant Rawls's moral psychology at the level of willingness to be fair, the stability story can still fail if citizens cannot reliably \emph{track} whether institutions and laws remain publicly justifiable over time. The rational/reasonable distinction is methodologically helpful here precisely because it blocks the easy dismissal---``your argument assumes people are bad''---while preserving the thought that defection pressures and incentive-driven drift are structurally present in politics. The critique targets the \emph{epistemic and cognitive} demands imposed by stability-through-public-justification, not Rawls's moral ideal of reasonable citizenship.

\section{Weithman Set-up: Stability as Equilibrium and Rawls's Late Concession}
Paul Weithman argues that Rawls's account of stability is best understood through the lens of \emph{equilibrium}. Rawls's question is not simply whether citizens comply, but whether just institutions can reproduce themselves by shaping citizens' motivations and beliefs in ways that support continued allegiance to those institutions. In Weithman's framing, Rawls's stability analysis (especially in \emph{A Theory of Justice}) is tightly connected to ideal theory and to the claim that justice and stability are linked through equilibrium-like relationships between institutions and citizen motivations.\parencite[p.~1188]{Weithman2024Rescuing}

The pivot that matters for Chapter~8 of this thesis begins with a concession in Rawls's own late work. In his reply to Habermas, Rawls allows that public political questions may be reasonably decidable not only by the most reasonable political conception of justice, but also by ``a reasonable family of such conceptions.''\parencite[p.~147]{rawls1995reply} He extends this line of thought in \emph{Public Reason Revisited}, where the content of public reason is explicitly given by ``a family of reasonable political conceptions of justice'' rather than by any single privileged view.\parencite[pp.~775--776]{rawls1997public} The significance of these passages is that Rawls himself opens the door to a set-valued account of what an overlapping consensus can sustain---a door Weithman subsequently walks through.

Weithman's 2023 analysis explains why this concession forces a conceptual shift. If a society can be effectively regulated by \emph{different} members of a set of Rawlsian liberal conceptions at different times, then the society is not ``well-ordered'' in the strict \emph{Theory of Justice} sense, which required effective regulation by a single conception.\parencite[p.~205]{Weithman2023FixedPoints} Rawls therefore needs a new idea: a \emph{well-ordered liberal society}, one that can be effectively regulated by any member of the liberal family and can transition among them over time.\parencite[p.~205]{Weithman2023FixedPoints} Weithman further argues that Rawls's concession ``would have to'' be expressed with a \emph{set-valued} equilibrium concept: a society may remain in equilibrium while transitioning within a set of acceptable justice points rather than converging on a single fixed point.\parencite[p.~206]{Weithman2023FixedPoints}

This Weithman set-up matters for the thesis in two ways. First, it offers Rawlsians a sophisticated rescue strategy: perhaps stability does not require convergence on a single conception, but only remaining within a narrow liberal neighborhood.\parencite[pp.~205--206]{Weithman2023FixedPoints} Second, it reframes the feasibility question that motivates the later critique. If stability is set-valued, then the citizen's justificatory task might appear less demanding---perhaps citizens need only track whether institutions remain within the liberal family. But this apparent relief depends on whether citizens and institutions can supply reliable, public, and temporally robust signals of what counts as within that family, and of when political drift has carried institutions beyond it. The remaining chapters examine whether Rawlsian political liberalism (even in its most charitable, Weithman-informed form) can plausibly meet those informational and cognitive demands.

\chapter{Feasibility Critique as Internal Critique}

Chapters~1--2 argued that Rawlsian political liberalism ties democratic legitimacy to a distinctive stability ideal: citizens are to regard coercive law as justifiable to them and to one another on the basis of public reasons, and that justificatory relationship is supposed to sustain a regime ``from one generation to the next.''\parencite[p.~1]{rawls1987overlapping} The present chapter clarifies what kind of critique the dissertation advances. The central challenge is not a familiar ``realist'' complaint that political philosophy should abandon idealization. Nor is it a complaint that citizens are too biased or selfish to deliberate. It is a feasibility critique aimed at a specific legitimating route within political liberalism: \emph{legitimacy through stability-by-public-justification}.

To position that critique, I first locate it within a broader tradition of feasibility objections (\S\,\ref{sec:feasibility-tradition}). I then argue that feasibility constraints are often \emph{internal} to Rawls's legitimacy ambitions, because stability is doing justificatory work in the theory rather than serving as an optional sociological add-on (\S\,\ref{sec:feasibility-types}). Finally, I bridge back to the thesis's two conclusions and preview the two ``invisible'' assumptions---\emph{tracking} and \emph{landscape stability}---that Chapter~4 will define more precisely (\S\,\ref{sec:bridge}--\S\,\ref{sec:preview}).

\section{The tradition of feasibility criticisms}\label{sec:feasibility-tradition}

Feasibility objections take many forms, but they share a basic structure: a normative proposal is said to fail, not (only) because it is false or unjustified, but because it cannot perform the role it is meant to perform for agents and institutions like ours.

A prominent contemporary entry point into these debates is the dispute over \emph{ideal theory}. Laura Valentini notes that ``ideal theory'' is often used loosely to describe theories constructed under false or idealized assumptions that make social reality ``significantly `simpler and better' than it actually is.''\parencite[p.~332]{valentini2009paradox} In the Rawlsian tradition, the most familiar idealizations are (i) full compliance and (ii) favorable background conditions---assumptions Rawls explicitly adopts in order to ``simplif[y] the problem of social justice.''\parencite[p.~332]{valentini2009paradox} Critics then press what Valentini calls the \emph{guidance critique}: principles formulated under idealized conditions will be ``ill-suited to guide action in the real world,'' and a sound normative theory is widely assumed to be action-guiding.\parencite[p.~333]{valentini2009paradox}

Yet Valentini also emphasizes an opposing pressure: giving up ideal theory does not look like a viable option, because judgments of justice and injustice (and plans for reform) seem to presuppose some ideal of what a just society would be like.\parencite[p.~333]{valentini2009paradox} That tension generates what she calls ``the paradox of ideal theory'': (a) any sound theory of justice is action-guiding; (b) any sound theory of justice is ideal; and (c) any ideal theory fails to be action-guiding.\parencite[p.~333]{valentini2009paradox} Whether and how that paradox can be resolved depends, among other things, on what kind of ``action-guidance'' a theory is supposed to provide, and what kinds of idealizations are permissible for that purpose.

Valentini's diagnosis also highlights a further fault line that matters for the present project. She distinguishes between \emph{fact-sensitive} approaches to justice---approaches that treat action-guidance as a necessary condition on a theory's practical adequacy---and \emph{fact-insensitive} approaches that deny that a theory of justice must be action-guiding in that way.\parencite[p.~333]{valentini2009paradox} In contemporary debates, G.~A.~Cohen is often taken to exemplify the latter stance.\parencite[p.~20]{cohen2008rescuing} This dissertation sides with Valentini's restriction in a limited sense: the object of critique is not the truth of ultimate moral principles in abstraction, but a \emph{legitimacy} ideal that presents itself as action-guiding for constitutional democracy.

David Estlund's \emph{Utopophobia} targets a different but overlapping strand of feasibility worry: the worry that political theory becomes objectionably ``utopian'' when it makes demands that seem unrealistic. Estlund characterizes \emph{utopophobia} as ``the unreasonable fear of the sin of utopianism,'' and he argues that this fear can marginalize valuable inquiry ``without demonstrating any defect'' in the standards proposed.\parencite[p.~6]{estlund2019utopophobia} The rhetorical charge that a theory is ``utopian,'' Estlund suggests, often smuggles in a substantive conclusion---that the theory's standards are false---by implying that they are ``too unrealistic,'' even though there are many distinct ways in which a theory might fail to be ``sufficiently realistic.''\parencite[p.~6]{estlund2019utopophobia}

Rawls's own work is frequently read as an attempt to occupy a middle ground between utopianism and concessionary realism. In \emph{The Law of Peoples}, Rawls defines political philosophy as \emph{realistically utopian} when it ``extends what are ordinarily thought of as the limits of practicable political possibility'' while remaining responsive to what is possible ``given the laws and tendencies of society.''\parencite[p.~11]{rawls1999peoples} That is: political philosophy can, and sometimes should, stretch our sense of what is politically possible, but it must still take seriously the background conditions under which political ideals can be institutionalized and sustained.

These debates set the stage for the present dissertation. The thesis does not claim that ideal theory is illegitimate as such. Nor does it claim that Rawls is guilty of the simplistic utopianism that critics often impute. Instead, it argues that Rawlsian political liberalism (and many of its refinements) relies on feasibility conditions that are easy to miss because they are embedded in the very way stability is supposed to generate legitimacy.

\section{Social, institutional, and conceptual feasibility}\label{sec:feasibility-types}

Feasibility criticisms are often treated as a single genre, but it is helpful to distinguish at least four families of constraints that can undermine a normative project's ability to perform its intended role.

\subsection{Motivational feasibility}

Some feasibility objections target motivation: can citizens and officials be expected (as a matter of moral psychology, incentives, or socialization) to comply with the demands of a theory? Rawls's own ideal theory assumes ``strict compliance,'' but his stability arguments---especially in political liberalism---are sensitive to whether citizens can be motivated by a sense of justice under conditions of pluralism.\parencite[p.~1]{rawls1987overlapping} The motivational critique is often the first reaction to public reason: perhaps citizens will not restrain themselves, perhaps they will not tolerate deep disagreement, perhaps they will not act from a shared political conception.

Estlund's discussion usefully complicates the picture by distinguishing two kinds of questions. On the one hand are \emph{concessive} requirements: what we ought to do given that we will not comply with other requirements (e.g., ``since we will not comply with socialist institutions in sufficient numbers, we ought not to build them'').\parencite[p.~6]{estlund2019utopophobia} On the other hand are \emph{nonconcessive} questions about what institutions we ought to build and comply with, even when we expect some noncompliance.\parencite[p.~6]{estlund2019utopophobia} The distinction matters because many feasibility objections move too quickly from motivational facts (noncompliance, corruption, polarization) to the conclusion that ideals are thereby discredited, rather than asking which parts of the normative problem are genuinely conditional on those motivational facts.

\subsection{Institutional feasibility}

A second family of objections targets institutional possibility: even if citizens were motivated, are there institutional arrangements that could realize the theory's requirements? In Rawls's own terms, the question is whether a political conception can be embodied in a basic structure that citizens can affirm and that can reproduce itself stably.\parencite[p.~1]{rawls1987overlapping} Institutional feasibility is not only about engineering: it includes questions about whether a society can sustain the legal forms, economic background conditions, and public political culture the theory presupposes.

\subsection{Epistemic and informational feasibility}
A third family of objections targets information: even if institutions could be designed and even if citizens were disposed to cooperate, do citizens (and officials) have access to the information required to satisfy the theory's justificatory conditions? Here Rawlsian political liberalism is distinctive. Public reason is not merely a compliance demand; it is a demand about \emph{justification} offered in the public political forum. Rawls frames that demand as part of ``the idea of democracy itself.''\parencite[p.~765]{rawls1997public}

Rawls identifies a specific informational burden built into this demand. A political conception requires not only publicly known principles but also ``publicly shared standards of inquiry and evidence''; without them, agreement on the conception is ``worthless---not an effective agreement at all.''\parencite[p.~8]{rawls1987overlapping} This is what Rawls elsewhere formalizes as the full publicity condition: the conception's complete justification must be available to all and acknowledged as the basis for political decisions, so that citizens can genuinely assess whether their institutions meet the standard of public reason. Samuel Freeman illuminates the practical upshot: a conception of justice is justified only if it can serve as ``a public basis in light of which citizens can justify to one another their common institutions,''\parencite[p.~11]{freeman2007burdens} and the publicity requirement extends to the reasoning that supports the principles, not merely to the principles themselves. Whatever one ultimately thinks about Rawls's political ideals, these demands impose a substantive informational burden: citizens must be able (in some sense) to access, understand, and apply the public basis of justification.

One reason epistemic feasibility objections are easily mischaracterized is that they can look like skeptical attacks on moral knowledge. Jonathan Quong argues that Rawls's requirement of ``epistemic restraint'' does \emph{not} presuppose skepticism about moral, religious, or philosophical truth.\parencite[p.~321]{quong2007without} Quong distinguishes a principle of restraint---do not appeal to contested comprehensive truths when deciding constitutional essentials\parencite[p.~321]{quong2007without}---from a principle of skepticism---doubting our ability to know such truths. Rawls cannot require skepticism, Quong argues, because skepticism about the good life is itself a controversial epistemic doctrine that many reasonable citizens reject; making legitimacy depend on it would undermine political liberalism's aim.\parencite[p.~322]{quong2007without} The central motivation for restraint is instead justificatory: the liberal state's basic structure must be justified in a way that all reasonable citizens can be expected to endorse.\parencite[p.~321]{quong2007without}

The thesis's feasibility critique is continuous with Quong's point. It does not argue that citizens lack access to moral truth, or that political liberalism secretly assumes skepticism. The worry is practical and diachronic: whether citizens can satisfy the informational and justificatory tasks that publicity and stability-through-reasons presuppose in complex, changing political environments.

The broader feasibility literature includes a large body of work emphasizing structural incentives toward political ignorance and the limits of epistemic outsourcing. Somin argues that, because the probability of any one citizen affecting electoral outcomes is tiny, the expected payoff from acquiring political knowledge is correspondingly low; \enquote{rational ignorance} is therefore endemic in mass democracy.\parencite[p.~77]{somin2013democracy}

Nor does shifting epistemic authority to experts solve the problem in general. Tetlock's study of expert political judgment finds \enquote{few signs that expertise translates into greater ability to make either `well calibrated' or `discriminating' forecasts}, even in the domains where experts are most confident.\parencite[p.~20]{tetlock2005expert}

Deliberative democrats respond by turning to institutional design, but Lafont insists that democratic legitimacy cannot be secured by epistemic \enquote{shortcuts} that bypass citizens' reflective participation in the name of better outcomes.\parencite[p.~3]{lafont2020democracy}

These literatures will matter later in the dissertation's diagnosis of the citizen-facing epistemic burden, but the present chapter's role is only to locate that burden within the feasibility tradition.

\subsection{Dynamic feasibility}

Finally, some feasibility objections target the dynamics of political life. A political conception might be motivationally and institutionally feasible in a snapshot sense while still failing to be feasible over time---because the environment in which justification occurs shifts, because institutional meanings drift, or because the mapping from public values to policy consequences changes.

Rawls's own stability aim is explicitly diachronic. He presents the point of an overlapping consensus as securing a shared basis of justification that can ``help ensure stability from one generation to the next.''\parencite[p.~1]{rawls1987overlapping} That diachronic orientation is part of what makes the feasibility question internal to Rawlsian legitimacy: the theory does not merely ask whether institutions would be just, but whether they can be stably justified to citizens across time.

\section{Why feasibility is an internal constraint here}\label{sec:bridge}

The foregoing typology explains why feasibility debates can easily become confused. Sometimes feasibility is an \emph{external} constraint on a normative ideal: one might think that moral truth and justice are independent of what people will do, and that feasibility is relevant only to strategy or implementation. That is one plausible way to understand certain forms of ideal theory.

But Rawlsian political liberalism makes feasibility harder to bracket, because stability is doing legitimating work. Rawls frames political philosophy's aim, in a constitutional democracy, as presenting a conception of justice that can provide ``a shared public basis for the justification of political and social institutions'' and ``help ensure stability from one generation to the next.''\parencite[p.~1]{rawls1987overlapping} The alternative is not simply ``less stable justice,'' but a qualitatively different basis of social unity: a mere \emph{modus vivendi} dependent on shifting interests and contingencies.\parencite[p.~1]{rawls1987overlapping} If legitimacy is tied to stability-through-public-justification, then facts about whether stability-through-justification is achievable are not merely implementation details; they bear on whether the legitimating route can succeed.

This internality becomes even clearer when we recall Rawls's characterization of public reason. Public reason, Rawls says, ``belongs to a conception of a well ordered constitutional democratic society,'' and its ``form and content'' are ``part of the idea of democracy itself.''\parencite[p.~765]{rawls1997public} These are not optional moral embellishments on top of coercive institutions. They are presented as conditions under which political power can be exercised in a way compatible with citizens' standing as free and equal.

The point is not that Rawls ignores feasibility. Indeed, Estlund notes that Rawls's ``concession to realism'' is ``no more than to limit his prescriptions to what is possible, however unlikely it might be,'' and Rawls at times describes his own project as ``realistic utopianism.''\parencite[p.~6]{estlund2019utopophobia} Rawls also explicitly defines political philosophy as realistically utopian: it extends the limits of practicable possibility while remaining tied to what is possible under the ``laws and tendencies of society.''\parencite[p.~11]{rawls1999peoples}

The thesis's distinctive claim is that Rawlsian political liberalism carries \emph{additional} feasibility commitments that are not always made explicit: commitments about what citizens and institutions must be able to do for stability to be ``for the right reasons,'' and commitments about how stable the justificatory environment must be for those tasks to be tractable. Recent work underscores how urgent the stability question remains in our political moment, even if Rawls's treatment of stability has historically received less attention than other parts of his project.\parencite[p.~1187]{Weithman2024Rescuing}

This is why the dissertation's critique is best understood as an \emph{internal} feasibility critique. It does not claim that Rawls's political values are unattractive, nor that the ideal of justification to free and equal citizens is misguided. It claims that the legitimating route that ties legitimacy to stability-through-public-justification depends on tacit assumptions about citizens' diachronic justificatory capacities and about the stability of the justificatory environment.

\section{Preview: tracking and landscape stability}\label{sec:preview}

The next chapter defines two assumptions that do hidden work in Rawlsian stability-through-justification.

\subsection{The tracking requirement (citizen-facing)}
If stability is to be ``for the right reasons,'' then citizens must be able to treat the political order as continuously justifiable to them on public grounds, not merely as something they happen to accept at a moment in time. In other words, they must be able to \emph{track} whether the justificatory status of coercive law is sustained as circumstances change. Rawls's own diachronic formulation---stability ``from one generation to the next''\parencite[p.~1]{rawls1987overlapping}---already implies that citizens' endorsement cannot be a one-off event.

Rawls's account of publicity makes the same point in a different register. The role of a mutually recognized political conception, he argues, is to provide a common point of view from which citizens can assess whether their institutions are just; political justification is addressed to others who disagree, and must be capable of serving as a shared basis for that assessment.\parencite[p.~6]{rawls1987overlapping} Freeman's discussion draws out the practical upshot: if a conception of justice must serve as a public basis in light of which citizens can justify their institutions to one another,\parencite[p.~11]{freeman2007burdens} then citizens must be able, at least at key moments, to recognize and apply that basis. Tracking is the diachronic extension of this thought: citizens must be able to recognize when the public basis still supports existing coercive outcomes, when it no longer does, and what kinds of revisions would restore justifiability.

\subsection{Landscape stability (world-facing)}

Even if citizens were perfectly reasonable and highly motivated, tracking requires an environment that is stable enough for tracking to succeed. The dissertation uses the metaphor of a ``justificatory landscape'' to describe the mapping from (i) public reasons and background facts, through (ii) institutional and policy mechanisms, to (iii) concrete coercive outcomes that citizens experience. If that mapping is highly nonstationary---if the meaning of constitutional concepts shifts quickly, or if technological and institutional changes rapidly alter the consequences of formally similar rules---then the justificatory task becomes harder even when the principles remain constant.

Chapter~4 makes these assumptions explicit and parameterizes them by timescale and granularity. The remainder of the dissertation then asks when the tracking requirement and landscape stability are likely to fail in modern, mediated democracies, and what becomes of Rawlsian legitimacy when they do.

\chapter{The Invisible Limits: Tracking and Landscape Stability}

Chapters~1--3 located the dissertation’s core worry within Rawlsian political liberalism’s own legitimating route. Rawls aims to present a political conception that can provide a shared public basis of justification for coercive institutions and thereby,\enquote {help ensure stability from one generation to the next}.\parencite[p.~1]{rawls1987overlapping} The critique developed so far is internal to that ambition: if legitimacy is supposed to be achieved through \emph{stability-by-public-justification}, then the relevant feasibility questions are not merely pragmatic questions about implementation. They are questions about whether the justificatory practice political liberalism requires can actually be sustained in the kinds of environments modern democracies inhabit.

This chapter makes explicit two assumptions that do hidden work in the stability-through-reasons story. The first is a \emph{tracking requirement}: citizens must be able to tell, over time, whether coercive outcomes remain publicly justifiable in terms of the political values they can reasonably endorse. The second is a \emph{landscape stability requirement}: the environment in which justification happens must be stable enough for that tracking task to be possible.

The chapter’s job is definitional and clarificatory. I do not yet argue that tracking fails or that landscapes are unstable (that is the work of Chapters~5--7). Instead, I (i) define the two requirements precisely enough to structure later debate; (ii) show why they are required by Rawlsian stability-for-the-right-reasons; and (iii) introduce two parameters---timescale (\(\tau\)) and granularity (\(g\))---that let us distinguish stronger from weaker versions of the requirements, and therefore stronger from weaker versions of the critique.

\section{The tracking requirement (citizen-facing)}\label{sec:tracking}

Rawls describes public reason as an idea that \enquote{belongs to a conception of a well ordered constitutional democratic society} and whose \enquote{form and content} are \enquote{part of the idea of democracy itself}.\parencite[p.~765]{rawls1997public} That formulation already suggests that public justification is not merely a technocratic procedure for officials. It is a constitutive part of democratic legitimacy: coercive law must be justifiable to citizens understood as free and equal.

Rawls also makes clear that the relevant justificatory task is temporally extended. The point of an overlapping consensus is not merely to secure momentary acquiescence, but to supply a shared basis of justification that can sustain a constitutional regime over time.\parencite[p.~1]{rawls1987overlapping} In his reply to Habermas, Rawls states the intended endpoint of this story with unusual bluntness: in a well-ordered constitutional democracy, when the basic structure is regulated by the most reasonable political conception (or a reasonable family of such conceptions) and is endorsed by an overlapping consensus, \enquote{from this follows stability for the right reasons}.\parencite[p.~147]{rawls1995reply}

The dissertation’s first hidden assumption is that this diachronic legitimacy ideal implicitly assigns citizens an epistemic task.

\subsection{A definition}

Call the political order at a time \(t\) \(\mathcal{O}_t\) (its constitutional structure, the basic structure more broadly, and the main coercive policies citizens confront).\footnote{This schematic notation is only a bookkeeping device. Chapter~6 introduces a toy model in prose; the formal material is optional and can be placed in an appendix, consistent with the revision plan’s emphasis on plain-English exposition.}

\begin{quote}
\textbf{Tracking (preliminary).} Citizens \emph{track} justificatory status over an interval \([t,t+\tau]\) when they are able, as circumstances change, to determine whether \(\mathcal{O}_t\) continues to be publicly justifiable to them (and to one another) by appeal to the political values and reasons that structure public reason.
\end{quote}

The definition is intentionally modest. It does not assume that citizens must know the truth of controversial moral or religious claims, nor that they must share comprehensive doctrines. Political liberalism explicitly rejects those requirements, and instead proposes that citizens must appeal to political values that others can reasonably endorse.\parencite[p.~765]{rawls1997public} Tracking is the diachronic extension of that same justificatory aspiration.

\subsection{Why Rawlsian stability implies tracking}

The link between stability-for-the-right-reasons and tracking can be drawn without adding new normative ideals to Rawls. If stability is to be secured \emph{by citizens’ reasons}---that is, by their continued ability to see coercive law as justified to them as free and equal---then their endorsement cannot be a single, once-and-for-all act. Rawls’s claim that a political conception should \enquote{help ensure stability from one generation to the next} signals a temporally extended standard.\parencite[p.~1]{rawls1987overlapping}

Public reason also gives this temporality a distinctive scope. Rawls insists that public reason applies to \enquote{fundamental political questions}---especially \enquote{constitutional essentials} and \enquote{matters of basic justice}.\parencite[p.~767]{rawls1997public} Because these questions recur as courts interpret constitutions, legislators revise basic legal frameworks, and social and technological conditions change, citizens’ standing as coauthors cannot be grounded in a one-time justificatory moment. The relevant question is whether citizens can continue to regard the exercise of political power as proper on the basis of reasons that others might also reasonably accept.

Rawls states this legitimacy condition in terms of reciprocity: political power is exercised properly only when we sincerely believe that our reasons are sufficient and when we can \enquote{reasonably think that other citizens might also reasonably accept those reasons}.\parencite[p.~771]{rawls1997public} That criterion is not just a snapshot requirement. It must be satisfiable across the sequence of occasions on which constitutional essentials and basic justice are at issue. Tracking, as defined here, is what citizens must be able to do if their allegiance is to remain allegiance \emph{for the right reasons} rather than a mere strategic accommodation to shifting power.

\subsection{What tracking is not}

It is important to rule out two misunderstandings.

First, tracking is not a demand that citizens continuously re-derive the whole political conception from first principles. Political liberalism is compatible with (and indeed relies on) a division of justificatory labor across institutions and time. What matters is that citizens can, at politically appropriate moments, assess whether coercive power is still exercised within a publicly justifiable framework.\parencite[p.~767]{rawls1997public}

Second, tracking is not the claim that citizens must converge on a single comprehensive worldview. Political liberalism’s whole point is that reasonable pluralism persists.\parencite[p.~765]{rawls1997public} Tracking asks whether, given persistent pluralism, citizens can still reliably locate political outcomes within the space of what public reason permits.

\section{Landscape stability (world-facing)}\label{sec:landscape}

Tracking is an epistemic task for citizens; landscape stability is a condition on the environment that makes that task tractable.

Rawls’s own writings do not use the metaphor of a \enquote{justificatory landscape}. But the metaphor is not gratuitous. Two Rawlsian themes invite it.

First, Rawls’s stability question is often interpreted (especially by Weithman) through equilibrium concepts. Weithman argues that Rawls’s key conceptual connections---\enquote{especially the connection between stability and equilibrium}---are laid out most clearly in \emph{A Theory of Justice} and are largely presupposed in \emph{Political Liberalism}.\parencite[p.~1188]{Weithman2024Rescuing} Equilibrium talk is naturally expressed in spatial terms: a system occupies a point or region and tends to remain there (or return there) under perturbation.

Second, the feasibility worry in this dissertation is explicitly dynamic. It concerns whether stability-by-public-justification can be sustained under conditions of social change and institutional adaptation. Those conditions are usefully illuminated by the complexity literature’s account of shifting payoff structures. Scott Page, for example, contrasts \emph{rugged landscapes}---where the relevant payoff function is fixed---with \emph{dancing landscapes}, where payoffs fluctuate because they depend on the evolving traits and strategies of other agents.\parencite[p.~94]{page2011diversity}

\subsection{A definition}

A \emph{justificatory landscape} is the mapping from public reasons and relevant facts to political outcomes and their justificatory status. Roughly:
\[
  (R_t, F_t) \xrightarrow{\;\text{institutions and rules}\; } O_t \xrightarrow{\;\text{public evaluation}\;} J_t.
\]
Here \(R_t\) is the set of public political reasons available in the culture at time \(t\); \(F_t\) is the relevant background of social facts (empirical regularities, technological constraints, institutional capacities); \(O_t\) is the vector of coercive outcomes produced by institutions; and \(J_t\) is the public-justificatory status of those outcomes (whether they can be justified, to citizens as free and equal, by appeal to public reasons).

With this in place:

\begin{quote}
\textbf{Landscape stability (preliminary).} A justificatory landscape is \emph{stable} over an interval \([t,t+\tau]\) at granularity \(g\) when the mapping from \((R,F)\) through institutions to outcomes and justificatory status does not change so rapidly or unpredictably that citizens cannot, at that granularity, track whether coercive outcomes remain publicly justifiable.
\end{quote}

Landscape stability is not the claim that societies must be static. Rawls expects a free society to contain permanent reasonable disagreement and ongoing political contestation.\parencite[p.~765]{rawls1997public} The claim is conditional: \emph{if} legitimacy is to be achieved by stability-through-public-justification, \emph{then} the justificatory environment must have enough stationarity for citizens’ justificatory judgments to remain answerable to something more than shifting contingencies.

\subsection{Why landscape stability is required}

Even perfectly reasonable citizens cannot track justificatory status if the mapping that determines that status is itself a moving target. Page’s model makes the relevant point in a stripped-down way. The rugged landscape model assumes that payoff does not change; but in real ecosystems, the value of a trait to one species depends on the traits of others, and that interdependence means that a fitness landscape \enquote{will fluctuate due to adaptations by other species}.\parencite[p.~94]{page2011diversity} Page calls these coupled landscapes \enquote{dancing landscapes}.\parencite[p.~94]{page2011diversity}

The political analogue is not biological fitness but public justification. Public reasons are applied through institutional mechanisms to produce coercive outcomes; citizens then evaluate those outcomes as justified or not. If changes in technology, institutional interpretation, or strategic adaptation systematically alter the consequences of formally similar rules, then the justificatory status of those rules may drift even when citizens’ professed political values remain constant. Under those conditions, the tracking task becomes more demanding because what must be tracked is not only whether citizens still affirm the same public values, but whether those values still support the same institutional outputs.

Holland’s description of complex adaptive systems sharpens the structural point. Governments and markets are not transparent machines; they are complex adaptive systems characterized by \enquote{intricate hierarchal arrangements of boundaries and signals}.\parencite[p.~1]{holland2012signals} Holland also emphasizes an epistemic limit: despite abundant data and description, we still know little about how to \enquote{steer} such systems, and widely applicable answers require studying the origins and coevolution of their signal and boundary hierarchies.\parencite[p.~1]{holland2012signals} If political institutions are complex adaptive systems of this kind, then changes in their internal signaling and boundary structures can change how public reasons translate into policy consequences---that is, they can change the justificatory landscape.

This is why the stability assumption is not a trivial simplification. The Rawlsian stability story is meant to secure more than episodic legitimacy; it aims at a durable public basis of justification.\parencite[p.~1]{rawls1987overlapping} That durability is achievable only if the reasons-to-outcomes mapping does not drift faster than citizens (and justificatory institutions) can reliably update their judgments.

\section{Timescale and granularity parameters}\label{sec:params}

So far, the requirements are stated in a way that invites an objection: perhaps they are too demanding. Rawls’s own account offers two natural ways to soften them. First, stability might be demanded only over certain timescales (not at every political moment). Second, public reason might operate at a coarse granularity (constitutional essentials and basic justice), not at the fine granularity of everyday policy.

To keep later debate honest, this section introduces two parameters.

\subsection{Timescale \(\tau\)}\label{sec:tau}

Rawls explicitly frames his stability aspiration in intergenerational terms. The aim is to offer a conception of justice that can provide a public basis of justification and \enquote{help ensure stability from one generation to the next}.\parencite[p.~1]{rawls1987overlapping} Weithman restates the same idea using standard philosophical language: Rawls takes stability to be the \enquote{propensity of a society, once well-ordered, to remain so over time}.\parencite[p.~198]{Weithman2023FixedPoints}

Yet Rawls also restricts public reason to questions that recur within political cycles. Public reason applies to \enquote{constitutional essentials} and \enquote{matters of basic justice}, including questions about what rights and liberties may reasonably be entrenched in a constitution and how those essentials are to be interpreted.\parencite[p.~767]{rawls1997public} Courts, legislatures, and democratic movements revisit those questions over timescales much shorter than a generation.

To capture both features, let \(\tau\) denote the interval over which stability-through-justification is demanded. A large \(\tau\) corresponds to Rawls’s intergenerational aspiration; a smaller \(\tau\) corresponds to the idea that legitimacy must be maintainable across shorter cycles of political and institutional change. Later chapters will argue that the tracking burden (and its failure modes) can look very different depending on \(\tau\).

\subsection{Granularity \(g\)}\label{sec:g}

Rawls’s doctrine of public reason is already a kind of coarse-graining. He insists that public reason has a definite structure and applies only to a limited subject matter: constitutional essentials and matters of basic justice.\parencite[p.~767]{rawls1997public} The doctrine is therefore not (and is not intended to be) a requirement that citizens track the justificatory status of every policy detail in real time.

We can represent this idea with a granularity parameter \(g\): the level of detail at which citizens and institutions individuate justificatory questions.

\begin{quote}
\textbf{Granularity.} At a fine-grained \(g\), citizens attempt to track justification at the level of particular policy mechanisms and empirical effects. At a coarse-grained \(g\), they track justification at the level of constitutional values, basic rights, and the family membership of political conceptions.
\end{quote}

Weithman’s interpretation of Rawls helps sharpen the strongest possible coarse-graining move. Rawls eventually conceded, Weithman argues, that an overlapping consensus is more likely to be on \enquote{a class of liberal conceptions} rather than a single conception such as justice as fairness.\parencite[p.~205]{Weithman2023FixedPoints} If a society can remain well-ordered while moving among the points picked out as just by different liberal conceptions, then stability cannot be identified with a single fixed point. Weithman concludes that Rawls’s concession requires a shift from a \emph{point-valued} equilibrium concept to a \emph{set-valued} one: a society may remain in equilibrium while transitioning within a set (or, in Weithman’s illustrative metric language, within a ball of acceptable liberal points).\parencite[pp.~205--206]{Weithman2023FixedPoints}

This is a powerful reply to an overly strong version of the tracking critique. If stability is set-valued at a coarse \(g\), then citizens do not need to track whether society occupies one uniquely correct justice point at all times. They might need only to track whether institutions remain within the acceptable liberal family.

\section{The granularity challenge stated cleanly}\label{sec:granularity-challenge}

The foregoing motivates the chapter’s final task: to state, in its strongest form, the Rawlsian challenge to the tracking/landscape critique.

\subsection{The challenge}

Rawlsians can plausibly argue that the critique has mis-specified what must be tracked. Public reason is not a demand for fine-grained policy monitoring, and Rawls explicitly limits its domain to constitutional essentials and basic justice.\parencite[p.~767]{rawls1997public} Moreover, the best contemporary reconstructions of Rawlsian stability (notably Weithman’s) suggest that stability can be understood as equilibrium in a \emph{set-valued} sense: stability might require remaining within a narrow neighborhood of reasonable liberal conceptions rather than converging on a single fixed point.\parencite[pp.~205--206]{Weithman2023FixedPoints}

If that is right, then the tracking requirement should be weakened along both dimensions. On the one hand, perhaps \(g\) is coarse: citizens need only track whether outcomes are compatible with liberal family membership. On the other hand, perhaps \(\tau\) is long: legitimacy is a matter of whether society remains within that family over generational time, not whether it remains there at every moment.

\subsection{What remains even after coarse-graining (preview only)}

The dissertation’s claim is not that this Rawlsian reply is confused or irrelevant. On the contrary, it is the most natural and most charitable way to reduce the burden on citizens while staying faithful to Rawls’s own restrictions.\parencite[p.~767]{rawls1997public} The claim is instead conditional and forward-looking:

\begin{enumerate}
  \item Even coarse-grained tracking requires reliable, publicly intelligible markers of what counts as \emph{within} the liberal family and of when drift has moved institutions beyond it. Weithman’s set-valued equilibrium is helpful precisely because it makes this boundary question explicit.\parencite[pp.~205--206]{Weithman2023FixedPoints}
  \item Landscape stability remains necessary even at coarse \(g\). If the mechanisms that translate constitutional values into concrete outcomes drift rapidly---through shifting interpretations, institutional redesign, or changing social facts---then citizens cannot reliably tell whether coercive outcomes remain within the liberal neighborhood that is supposed to secure legitimacy.
\end{enumerate}

Those are not yet arguments; they are signposts. Chapters~5--7 will argue that modern, mediated democracies often produce exactly the kind of drift that makes these two problems acute. For now, Chapter~4’s result is conceptual: Rawlsian stability-through-public-justification presupposes (i) a citizen-facing capacity to track justificatory status over time and (ii) a sufficiently stable justificatory landscape for that tracking to be feasible. Once those assumptions are explicit and parameterized by \(\tau\) and \(g\), the later chapters can ask, with greater precision, when Rawlsian legitimacy can plausibly be achieved and when it cannot.

\chapter{Why These Assumptions Fail in Mediated Politics}

Chapter~4 made explicit two requirements that do hidden work in Rawls’s stability-through-reasons story: (i) a \emph{citizen-facing} requirement that citizens be able to track justificatory status over time, and (ii) a \emph{world-facing} requirement that the justificatory landscape be stable enough for that tracking to succeed. This chapter argues that both requirements are empirically fragile in the settings where modern constitutional democracies actually operate: highly mediated political environments in which citizens encounter political facts, opponents, and justificatory narratives primarily through intermediaries rather than through direct, face-to-face public reasoning.

The claim is not that mediation is a regrettable accident that a better civics curriculum could simply wash away. Rawls’s own account distinguishes the \enquote{public political forum} from the broader \enquote{background culture} and makes clear that the idea of public reason does not apply to the background culture, including \enquote{media of any kind}.\parencite[pp.~767--768]{rawls1997public} That division of labor is meant to protect liberal legitimacy from the impossible task of purifying an entire society’s cultural life. But it also generates a feasibility dependence: if citizens’ beliefs and affective orientations are shaped outside the public forum, then stability-by-public-justification depends on whether the background culture supplies a sufficiently stable and mutually intelligible environment for citizens to \emph{track} what is publicly justifiable.

This chapter provides evidence for a focused diagnosis. Contemporary mediated politics often destabilizes (a) the empirical predicates on which justifications rely and (b) the affective and interpretive conditions under which citizens treat one another as eligible recipients of public justification. When those background conditions fail, citizens can remain good-faith reasoners and still be unable to track whether coercive outcomes remain publicly justifiable.

\section*{Evidence anchors (and what they do in this chapter)}

Because this chapter is evidence-sensitive, it is useful to state explicitly how the primary empirical anchors function in the argument.

\begin{itemize}
  \item \textbf{Mediated polarization \& the erosion of cross-cutting understanding.} Levendusky and Malhotra show experimentally that media coverage depicting a polarized electorate increases perceived polarization\parencite[p.~290]{levenduskyMalhotra2016polarization} and increases affective polarization toward the opposing party.\parencite[p.~292]{levenduskyMalhotra2016polarization} These effects show that a mediated narrative about \emph{how divided we are} can change both citizens’ descriptive beliefs and their evaluative stances toward opponents. Mutz complements this by showing what is lost when contact across political difference is scarce or erodes: cross-cutting exposure can increase awareness of legitimate rationales for opposing views and contribute to political tolerance,\parencite[p.~118]{mutz2002crosscutting} yet cross-cutting exposure can be \enquote{disappointingly infrequent} and is plausibly threatened by processes like residential balkanization.\parencite[pp.~111, 122]{mutz2002crosscutting} Combined, these sources support the chapter’s central claim that mediated polarization can undermine the citizen-facing task of tracking justification by thinning ordinary social mechanisms for learning what opponents’ reasons look like.
  \item \textbf{Cascades, diffusion, and polarization dynamics.} Bikhchandani, Hirshleifer, and Welch provide a canonical mechanism by which small informational shocks become large behavioral convergence: an \emph{informational cascade} occurs when it becomes optimal to follow predecessors \enquote{without regard to} one’s own information.\parencite[p.~992]{bikhchandaniHirshleiferWelch1992cascades} They emphasize the \emph{fragility} of cascades, including their susceptibility to abrupt change.\parencite[p.~996]{bikhchandaniHirshleiferWelch1992cascades} Vosoughi, Roy, and Aral show that, in a contemporary social-media setting, false news diffuses \enquote{farther, faster, deeper, and more broadly} than true news, especially in politics.\parencite[p.~1146]{vosoughiRoyAral2018falseNews} Sunstein adds a deliberative analogue: group discussion often produces \emph{group polarization}, shifting people toward more extreme positions in the direction of their initial tendencies.\parencite[p.~176]{sunstein2002groupPolarization} Combined, these sources anchor the chapter’s claim that the informational and social dynamics of mediated politics can make the justificatory landscape shift quickly, unpredictably, and sometimes discontinuously.
\end{itemize}

\section{Unstable justification}\label{sec:ch5-unstable}

The stability-through-reasons ideal presupposes more than citizens’ sincerity. It presupposes that citizens can form reasonably accurate beliefs about the social facts that bear on constitutional essentials and basic justice and that the salience of political values does not shift so quickly that yesterday’s public justification becomes today’s opaque slogan. In a mediated environment, both presuppositions can fail.

This section isolates two channels of instability, corresponding to the chapter outline: shifting empirical regularities (which change how institutional rules actually operate) and shifting moral salience (which changes what citizens are prepared to treat as a decisive public reason).

\subsection{Shifting empirical regularities affecting policy consequences}\label{sec:ch5-empirical}

In Rawls’s framework, citizens are meant to justify coercive political outcomes by reference to political values applied to a shared set of social facts.\parencite[p.~765]{rawls1997public} But the factual environment relevant to politics is not merely a static backdrop. In a mediated society, the epistemic ecology through which facts are encountered is itself dynamic, and that dynamism can alter the practical meaning of institutions that appear stable at the level of formal description.

A useful starting point is the mechanism of informational cascades. Bikhchandani, Hirshleifer, and Welch define an informational cascade as a situation in which, after observing the actions of those ahead, it is optimal for an agent to follow the predecessors’ behavior \enquote{without regard to his own information}.\parencite[p.~992]{bikhchandaniHirshleiferWelch1992cascades} The appeal of this mechanism is not that it imputes irrationality. On the contrary, cascades can be individually rational responses to environments in which other people’s actions are treated as evidence. But cascades change the relationship between evidence and belief in two politically important ways.

First, they provide a general explanation of rapid convergence. If visible social signals are treated as informative, then early movers (or salient actors) can have disproportionate influence over what becomes publicly accepted. This matters for public justification because political argument is not only about first-order values. It is also about empirical beliefs: whether a policy is effective, whether a group is threatening, whether an institutional rule is being exploited, whether others are complying. When those empirical beliefs are formed under cascade conditions, the public factual basis that citizens use to evaluate policy can shift quickly even if the underlying policy problem has not changed.

Second, cascades create fragility. The same mechanism that produces swift convergence can produce swift reversal. Bikhchandani, Hirshleifer, and Welch explicitly emphasize the \enquote{fragility of cascades} and their sensitivity to certain kinds of shocks.\parencite[p.~996]{bikhchandaniHirshleiferWelch1992cascades} In political settings, this fragility matters because justificatory status often depends on public beliefs about risks, prevalence, and consequences. When those beliefs are highly path-dependent and easily destabilized, the mapping from political values to justified policy can become unstable, even if the values remain constant.

Recent work on online diffusion suggests that cascade-like dynamics are not merely theoretical. Vosoughi, Roy, and Aral study verified true and false stories on Twitter from 2006 to 2017 and find that false news diffuses significantly \enquote{farther, faster, deeper, and more broadly} than true news; they report that the effect is especially pronounced for political news.\parencite[p.~1146]{vosoughiRoyAral2018falseNews} Their headline result is empirical and comparative: the informational environment is not neutral with respect to accuracy. If false political news systematically travels farther and faster than true political news, then the public epistemic environment can drift toward distorted empirical premises at precisely the moments when high-stakes political judgment is demanded.

These observations bear directly on landscape stability. In Chapter~4, landscape stability was defined as a property of the mapping from reasons and facts to institutional outcomes and their justificatory status. A mediated environment with cascades and asymmetric diffusion undermines that stability in at least two ways. It can generate rapid changes in what citizens take to be the relevant facts, and it can generate rapid changes in the perceived consequences of institutional rules (since perceived consequences are themselves mediated by what people believe and expect others to do).

This point does not require the extreme conclusion that truth is impossible in politics. It requires only a weaker and more plausible claim: that the empirical predicates relevant to public justification can fluctuate in ways that outpace ordinary citizens’ ability to update responsibly. When the factual environment dances in that way, tracking justificatory status becomes difficult even for reasonable citizens, because the target they are supposed to track is not fixed.

\subsection{Shifting moral salience that reweights political values}\label{sec:ch5-salience}

Political liberalism is not merely a theory of social facts; it is a theory of \emph{public justification}. It therefore relies on citizens’ willingness to treat certain political values as reasons that others can reasonably accept. But mediated environments can alter the affective and identity-laden background against which citizens weigh those political values.

Levendusky and Malhotra provide a clean case of salience shifting through mediated description. Their central question is not whether the public is polarized in some metaphysical sense, but what happens when the media portrays the electorate as polarized. In their experimental design, the \enquote{polarized} article increases perceived polarization between Democrats and Republicans.\parencite[p.~290]{levenduskyMalhotra2016polarization} More importantly for legitimacy-through-justification, the polarized coverage increases affective polarization: subjects exposed to polarized coverage rate opposing partisans lower on feeling thermometers and report more negative reactions to the out-party.\parencite[p.~292]{levenduskyMalhotra2016polarization}

For Rawls, the reciprocity condition is meant to make political power acceptable among free and equal citizens: one must be able to offer reasons one can \enquote{reasonably think} others might accept.\parencite[p.~771]{rawls1997public} The feasibility of that stance is not purely cognitive; it is also affective and social. When mediated narratives intensify dislike of the out-party, they can make citizens less disposed to interpret opponents charitably, less willing to treat opponents as reasonable, and more inclined to treat political conflict as evidence of bad faith. In other words, the same society can satisfy the formal criteria of constitutional democracy while losing some of the informal interpretive conditions that make public justification psychologically and sociologically sustainable.

Sunstein’s account of group polarization adds a complementary social dynamic. He argues that group polarization is a striking statistical regularity: after deliberation, members of a group tend to move toward \enquote{a more extreme point in the direction indicated by the members’ predeliberation tendencies}.\parencite[p.~176]{sunstein2002groupPolarization} When people repeatedly engage in discussion with like-minded others without sustained exposure to competing views, Sunstein argues, extreme movements become more likely.\parencite[p.~176]{sunstein2002groupPolarization} This is directly relevant to modern mediation: mediated environments often sort citizens into interaction patterns that are not representative of the polity as a whole. Under those conditions, deliberation can magnify rather than attenuate moral and political divergence.

Mutz’s study of cross-cutting networks shows what is at stake when cross-cutting contact is scarce. Mutz frames her inquiry by noting that exposure to conflicting political viewpoints is widely assumed to benefit democratic citizens.\parencite[p.~111]{mutz2002crosscutting} Her results speak to a specific epistemic and moral function of cross-cutting exposure: it can increase citizens’ awareness of \enquote{legitimate rationales for opposing views} and, through cognitive and affective channels, contribute to political tolerance.\parencite[p.~118]{mutz2002crosscutting} Yet Mutz also notes that cross-cutting exposure can be \enquote{disappointingly infrequent}.\parencite[p.~111]{mutz2002crosscutting} And she emphasizes that trends like \enquote{residential balkanization} plausibly translate into less communication across political difference, reducing opportunities for citizens to learn rationales for opponents.\parencite[p.~122]{mutz2002crosscutting}

Together, these findings illustrate why the landscape stability requirement is not merely about \enquote{facts}. It is also about what citizens treat as salient reasons. In a social setting where polarization dynamics intensify negative affect and reduce cross-cutting understanding, citizens’ weighting of political values can shift. For example, the political value of reciprocity (treating opponents’ reasons as reasons that might be reasonable) can lose motivational grip as out-party dislike rises. Or the political value of civility can become interpreted as weakness in a polarized setting. Those shifts are not mere psychological noise; they change which public reasons citizens will count as decisive, and therefore change the justificatory status of policies without any formal change in the underlying political conception.

\section{Tracking in a mediated political sphere}\label{sec:ch5-tracking-mediated}

The preceding section described why public justification can become unstable. This section connects that instability to the tracking requirement.

Rawls restricts public reason to \enquote{constitutional essentials} and \enquote{matters of basic justice} and stresses that it governs discourse in the public political forum (judges, officials, candidates), not the broader background culture.\parencite[p.~767]{rawls1997public} He also says explicitly that public reason does not apply to \enquote{media of any kind}.\parencite[p.~768]{rawls1997public} This is often presented as a strength of political liberalism: it avoids treating liberal legitimacy as a demand to control or censor the entire cultural sphere. But the same division of labor makes a feasibility dependence unavoidable. Citizens do not form their political beliefs and affective orientations only in the public forum; they form them primarily in the background culture. If the background culture systematically destabilizes the reasons-to-judgment mapping, then citizens cannot be assumed to track justificatory status reliably, even if they are willing to reason fairly.

Tracking requires at least two minimal capacities. Citizens must have some access to relevant social facts and policy consequences, and they must have some stable ability to see opponents’ reasons as eligible inputs into public justification. Both capacities can be impaired by mediated dynamics highlighted above.

First, the polarization evidence suggests that citizens’ descriptive beliefs about \emph{what kind of society they inhabit} are sensitive to journalistic framing. Levendusky and Malhotra show that exposure to polarized coverage increases perceptions that the electorate is polarized.\parencite[p.~290]{levenduskyMalhotra2016polarization} That matters for tracking because public justification is relational. Whether a political value plausibly supports a coercive policy often depends on beliefs about what others will do, what others will accept, and what threats are present. If citizens systematically update those background beliefs in response to media portrayals of division, then the justificatory landscape is not a stable object to be tracked.

Second, the same media portrayals can change affective postures. Levendusky and Malhotra find that exposure to polarized coverage increases dislike of the opposing party.\parencite[p.~292]{levenduskyMalhotra2016polarization} In Rawlsian terms, this matters because the exercise of political power is supposed to be justifiable to citizens understood as free and equal, and the reciprocity condition requires that we offer reasons we can reasonably believe others might accept.\parencite[p.~771]{rawls1997public} Affective polarization makes that reciprocity stance harder to sustain; it increases the probability that citizens will interpret opponents as unreasonable, and therefore treat opponents’ reasons as illegitimate by default. Even if citizens remain committed to liberal values, the interpretive environment needed for those values to do legitimating work can degrade.

Third, Mutz shows that cross-cutting exposure has a distinctive justificatory function. One reason public justification is possible, in practice, is that citizens have everyday routes for learning what opponents’ rationales look like. Mutz finds that exposure to dissonant views contributes to awareness of rationales for opposing views and can increase political tolerance.\parencite[p.~118]{mutz2002crosscutting} If cross-cutting exposure is rare\parencite[p.~111]{mutz2002crosscutting} and threatened by patterns of social sorting,\parencite[p.~122]{mutz2002crosscutting} then citizens’ ability to track justification can fail even at coarse granularity. They may still endorse constitutional values, but lack practical knowledge of how those values are interpreted by others, or of what compromises others could accept. Tracking requires not only stable values but stable \emph{public intelligibility} of reasons.

Finally, the second evidence-anchor set explains why these pressures can be self-reinforcing. Cascades show how individuals can rationally ignore private information once a social signal becomes salient.\parencite[p.~992]{bikhchandaniHirshleiferWelch1992cascades} The fragility of cascades suggests that even small shocks can produce abrupt shifts in what becomes publicly accepted.\parencite[p.~996]{bikhchandaniHirshleiferWelch1992cascades} Vosoughi, Roy, and Aral show that in an online environment falsehood spreads faster and farther than truth.\parencite[p.~1146]{vosoughiRoyAral2018falseNews} Sunstein shows that group deliberation often shifts people toward more extreme positions, especially when like-minded groups interact repeatedly without sustained exposure to opposing views.\parencite[p.~176]{sunstein2002groupPolarization} In combination, these results describe a background culture that can create rapid changes in public belief and affect, and can convert those changes into more entrenched polarization through group interaction.

None of this proves that Rawlsian legitimacy is \emph{never} achievable. It does show, however, that legitimacy-through-stability-by-public-justification depends on background conditions that modern mediated politics often fails to supply. The Rawlsian division between public forum and background culture is not an afterthought; it is the location where the feasibility problem becomes visible. If public reason does not govern the media,\parencite[p.~768]{rawls1997public} then Rawlsian stability depends on whether media dynamics are sufficiently benign that citizens can still meet the tracking requirement. The empirical evidence reviewed here suggests that this is not a safe assumption.

\section{Interim conclusion}\label{sec:ch5-interim}

If stability-for-the-right-reasons requires citizens to track justificatory status over time, then it matters whether the background culture supplies stable, public, and mutually intelligible cues about what reasons others could accept. Rawls’s own insistence that public reason does not apply to \enquote{media of any kind} makes this dependence explicit.\parencite[p.~768]{rawls1997public}

The evidence reviewed here identifies two connected pressures in mediated politics. On the one hand, information dynamics can produce rapid and fragile convergence in belief (cascades) and systematically asymmetric diffusion of falsehood online.\parencite[pp.~992, 996]{bikhchandaniHirshleiferWelch1992cascades}\parencite[p.~1146]{vosoughiRoyAral2018falseNews} On the other hand, media depictions of polarization can increase perceived and affective polarization,\parencite[pp.~290, 292]{levenduskyMalhotra2016polarization} and interaction among like-minded groups can reliably intensify extremity.\parencite[p.~176]{sunstein2002groupPolarization} When these dynamics thin cross-cutting social contact, citizens lose ordinary mechanisms for learning legitimate rationales for opposing views.\parencite[pp.~111, 118, 122]{mutz2002crosscutting}

The upshot is not that citizens are irrational or ill-motivated. It is that the social machinery through which citizens encounter one another and the world can make the justificatory landscape move on timescales shorter than the stability ideal can tolerate. The next chapters show how these mediated instabilities interact with bounded cognition and with the dynamics of complex political environments, tightening the claim that Rawlsian stability-through-reasons is feasibility-sensitive in ways political liberalism too often leaves implicit.

\chapter{JTC and Cognitive Burden}
\label{ch:jtc-cog-burden}

Rawlsian political liberalism ties legitimacy to a distinctive kind of stability: a regime is stable \emph{for the right reasons} when citizens comply with its basic structure because they can endorse it, not merely because they are coerced or strategically accommodated. As Paul Weithman emphasizes, a large share of Rawls's work on justice is devoted to the question of whether just institutions would be \emph{stably} just. \parencite[p.~1188]{Weithman2024Rescuing}

Chapters~4--5 argued that this stability ambition implicitly presupposes a \emph{tracking} capacity: citizens must be able, over time and across changing conditions, to recognize when coercive political power continues to be exercised on terms they can reasonably accept. This chapter tightens that diagnosis by asking a more basic question: \emph{What does that tracking task require of citizens as cognitive agents?}

The central claim is not that citizens are irrational, selfish, or apathetic. The claim is structural: even citizens who are sincerely trying to act as Rawlsian citizens face severe epistemic and computational constraints. Contemporary neuroscience and cognitive science increasingly model perception and cognition as forms of \emph{predictive inference} that trade accuracy for tractability, and that compress high-dimensional information into sparse, low-dimensional representations. When the justificatory environment drifts---when the consequences, salience, and institutional meaning of policies shift over time---those constraints generate a standing pressure toward mis-tracking. This pressure does not disappear even if we accept the Rawlsian narrowing of public reason to constitutional essentials and basic justice.

\section{JTC and Rawlsian legitimacy}
\label{sec:jtc-legitimacy}

A useful way to capture Rawls's legitimacy aspiration is what I have been calling the \emph{justification-to-citizens} requirement (JTC): coercive political power is legitimate only if it is justified in terms that citizens can reasonably accept as free and equal persons.

Rawls's own formulation of the public-reason version of this idea is explicitly interpersonal. In exercising political power, we should not merely \emph{have} reasons; we must be able to offer reasons that satisfy a reciprocity constraint. Rawls writes that our exercise of political power is proper only when we sincerely believe the reasons we would offer for political action are sufficient \emph{and} we reasonably think other citizens might also reasonably accept those reasons. \parencite[pp.~446--447]{rawls2005political}

This is not merely a constraint on public officials. Rawls connects the idea of public reason to a \emph{duty of civility} that is meant to shape political conduct and self-understanding. In a representative democracy, citizens exercise political power by voting, and they should therefore \enquote{think of themselves as ideal legislators} asking what statutes and policies can be supported by reasons that satisfy reciprocity. They should also be disposed to repudiate officials who violate public reason. \parencite[p.~444]{rawls2005political} The point, in Rawls's framework, is that citizens do not merely \emph{receive} justification; they participate (directly or indirectly) in its ongoing maintenance.

Rawls also insists that public reason has a definite subject matter. Political conceptions of justice should be \emph{complete} in a practical sense: they should specify values and principles that can answer \enquote{all or nearly all} questions involving \emph{constitutional essentials} and \emph{matters of basic justice}. \parencite[pp.~453--454]{rawls2005political} This narrowing is often invoked as a reply to feasibility worries: public reason is not supposed to govern every policy detail, so citizens are not asked to track everything.

But even in this narrowed form, Rawls's account has a demanding cognitive upshot. Public reason is meant to be publicly shareable, supported by citizens' allegiance to it, and sustained by a vibrant political culture; Rawls warns that when citizens \enquote{no longer see the point} of affirming the ideal of public reason, they may \enquote{fall into bitterness and resentment}. \parencite[p.~485]{rawls2005political} The question, then, is whether the cognitive task that JTC assigns to citizens is one that finite agents can realistically sustain under drift.

\section{Predictive brains, compressed worlds}

The core idea of predictive-processing approaches is that perception and action are forms of inference under uncertainty: organisms continuously generate expectations about their sensory stream and then revise those expectations in light of mismatch. One canonical implementation is \emph{predictive coding}, in which backward (top--down) connections convey predictions while forward (bottom--up) connections convey only the residual \emph{prediction error}.\parencite[p.~816]{friston2005theory} Clark captures the intuitive picture: the system uses \enquote{what you already know} to \emph{explain away} the incoming signal, updating the model when prediction error persists.\parencite[p.~183]{clark2013whatever}

The motivation for this architecture is computational. The agent faces a \emph{recognition} problem: given a stream of sensory data, infer the hidden causes that produced it. In realistic settings, the mapping from causes to sensations is nonlinear and context-sensitive, so there is generally \enquote{no analytic solution} for inverting the generative process.\parencite[p.~816]{friston2005theory} Friston proposes to treat inference, learning, and action as a single optimization problem: minimize \emph{variational free energy}, an objective function that upper-bounds surprise and trades representational \emph{accuracy} against model \emph{complexity}.\parencite[pp.~127--128]{friston2010freeenergy}

But any such solution is also a compression scheme. The world presents a high-dimensional, noisy signal; internal models must be lossy enough to be tractable while stable enough to guide action. Ganguli and Sompolinsky emphasize that neural systems routinely construct compressed representations that preserve task-relevant information while discarding detail.\parencite[pp.~491--492]{ganguli2012compressed} Work on sparse distributed representations makes the same point from a coding perspective: robustness and capacity depend on extreme sparsity, which is itself a kind of structural coarse-graining.\parencite[p.~1]{ahmadHawkins2016sdr}

The political upshot is not that citizens are \emph{irrational}. It is that even ideally motivated agents will rely on compressive, prediction-driven models when the environment is complex, partially observable, and strategically mediated---exactly the conditions under which citizens encounter the justificatory currency of public reason.
\section{Drift as a defeater for justificatory belief}
\label{sec:drift-defeater}

Call \emph{drift} any diachronic change in the mapping from policies and institutional rules to their salient consequences and meanings---changes in technology, economic structure, administrative practice, interpretive doctrine, or social norms that alter what a policy actually does \emph{and} what it is taken to signify. Drift matters for Rawlsian legitimacy because JTC is not merely a one-shot requirement. It is a continuing constraint on political power: citizens must be able, over time, to see coercion as exercised on terms they can reasonably accept. \parencite[pp.~446--447]{rawls2005political}

Predictive-processing frameworks suggest a distinctive way in which drift can defeat justificatory belief. Because the system is organized around prediction and error correction, updates are driven by the detection of \emph{mismatch}. When mismatch is small or ambiguous, the system has multiple ways to reduce it: it can revise the internal model, discount the signal as noise, or act to change the incoming sensory stream. Friston explicitly treats prediction errors as signals that \enquote{try to suppress themselves} through action, coupling perception and action in the control of error. \parencite[p.~134]{friston2010freeenergy} Clark likewise emphasizes that prediction-error minimization is not confined to perception but extends to action and learning. \parencite[p.~190]{clark2013whatever}

In politics, however, much of what matters for justification is not a sharp, unambiguous error signal. The relationship between a constitutional principle and its institutional implementation is typically mediated, interpretive, and contested. That makes drift easy to miss: slow shifts can accumulate under the threshold of perceived mismatch, especially when citizens' models compress the environment and favor simpler explanations over complex ones. \parencite[p.~128]{friston2010freeenergy}\parencite[pp.~491--492]{ganguli2012compressed} When the mismatch finally becomes salient, citizens may experience it not as a prompt for fine-grained revision but as a global shock that undermines trust in the justificatory scheme as a whole---the kind of affective and motivational failure Rawls himself worries about when citizens no longer see the point of the public reason ideal. \parencite[p.~485]{rawls2005political}

A helpful metaphor here comes from Scott Page's account of \enquote{dancing landscapes}. Page emphasizes that when payoffs depend on the adaptations of others, the \enquote{fitness landscape} fluctuates; he calls the resulting coupled landscapes \enquote{dancing landscapes} and treats them as a model of co-evolutionary change. \parencite[p.~94]{page2011diversity} Without importing the full complexity apparatus (reserved for the next chapter), the basic lesson is enough: a citizen can be locally rational and sincerely motivated, yet still fail to track a moving target.

The result is a defeater structure: citizens may continue to endorse coercive laws on the basis of reasons that were once publicly shareable and reciprocity-satisfying, even after drift has changed the relevant mapping from those laws to their actual consequences and public meaning. Under JTC, that is a legitimacy problem \emph{internal} to the ideal: coercion is supposed to remain justifiable to citizens \emph{as they are}, in a society whose facts and interpretations evolve.

\section{A toy model of justificatory tracking}
\label{sec:toy-model}

The point of a toy model is not to reduce political legitimacy to neuroscience or to a single equation. It is to make explicit the informational structure of the tracking burden.

Consider three layers:

\begin{enumerate}[leftmargin=*, itemsep=0.4em]
    \item \textbf{A world-layer:} a set of institutional and empirical facts at time $t$ that determine how a policy actually functions and what it produces.
    \item \textbf{A justification-layer:} a public-reason standard that, given those facts, classifies certain uses of coercion as acceptable (reciprocity-satisfying) and others as not.
    \item \textbf{A citizen-layer:} citizens' internal representations of the world-layer and the justification-layer, used to decide whether a policy remains acceptable.
\end{enumerate}

JTC requires that the citizen-layer stay aligned with the justification-layer over time. Rawls's duty-of-civility picture makes the alignment responsibility explicit: citizens are to assess, as if they were legislators, whether statutes can be supported by public reasons satisfying reciprocity, and they are to repudiate violations. \parencite[p.~444]{rawls2005political}

Now add drift: the world-layer changes over time, so the mapping from policy to consequences and meanings changes. The citizen must update. But predictive-processing models describe updates as driven by prediction error and governed by an accuracy/complexity trade-off. \parencite[p.~128]{friston2010freeenergy} And because the underlying information is high-dimensional, citizens will compress it into sparse, low-dimensional summaries. \parencite[pp.~491--492]{ganguli2012compressed} The toy model therefore predicts a familiar pattern: citizens can maintain approximate alignment in stable periods, but alignment will degrade under sustained drift unless there are strong, reliable, publicly accessible error signals that trigger revision.

Appendix~\ref{app:formal-model} sketches one way to formalize this idea without turning the main line of argument into a technical exercise.

\section{The granularity reply and what survives it}
\label{sec:granularity-reply}

A natural Rawlsian reply is that the tracking task has been overstated. Public reason, after all, does not demand that citizens monitor every policy detail. It targets constitutional essentials and basic justice, and political conceptions are meant to supply principles that can answer \emph{those} questions. \parencite[pp.~453--454]{rawls2005political} Moreover, the content of public reason is not fixed by a single doctrine but is given by a \emph{family} of reasonable liberal political conceptions satisfying reciprocity. \parencite[p.~485]{rawls2005political} If legitimacy only requires staying within that family, the cognitive burden appears much lighter.

This reply has real force. It motivates the most charitable version of Rawlsian feasibility: citizens can be coarse-grained trackers who mainly need to know whether their institutions remain within a broadly liberal range.

However, two points survive even after we grant the reply.

First, Rawls's own duty-of-civility account still assigns citizens an active monitoring role: they are to view themselves as legislators, ask which statutes can be supported by reciprocity-satisfying reasons, and repudiate officials who violate public reason. \parencite[p.~444]{rawls2005political} That role presupposes more than a vague sense that \enquote{we are still liberal}. It presupposes that citizens can discriminate (at least roughly) between coercion that remains publicly justifiable and coercion that has drifted into forms others cannot reasonably accept.

Second, even coarse-grained membership judgments are themselves vulnerable to drift. Rawls explicitly frames the content of public reason as a \enquote{family} of reasonable liberal political conceptions, not a single canonical doctrine.\parencite[p.~450]{rawls2005political} The boundary of such a family is learned, interpreted, and socially maintained. Predictive-processing accounts explain why agents tend to prefer simpler models and compress information under uncertainty, trading complexity for tractability.\parencite[p.~128]{friston2010freeenergy}\parencite[pp.~491--492]{ganguli2012compressed} When institutional meaning evolves---for example, when a practice retains its legal label while its real-world function shifts---citizens' coarse categories can lag behind reality.

The conclusion is not that Rawlsian public reason is a worthless ideal. The conclusion is that, once we model citizens as finite predictive agents operating under informational compression and drift, JTC imposes a demanding diachronic tracking task whose feasibility is far from guaranteed. The next chapter widens the lens by turning from cognitive constraints to environmental dynamics: even if citizens could update perfectly, a sufficiently \enquote{dancing} justificatory landscape would defeat stability-through-reasons as a durable political achievement.

\chapter{Complexity and Political Drift}

\noindent This chapter extends the thesis’s argument about the limits of public reason by adding a structural constraint that is neither purely moral nor purely cognitive: the complexity of the social world in which public justification must operate. Even when citizens share a political conception and even when deliberators act in good faith, the justificatory landscape in an open society is not static. It changes as institutions, incentives, technologies, and social identities co-evolve. In Scott Page’s terms, political justification is not carried out on a fixed landscape; it often occurs on \emph{dancing landscapes} whose peaks shift as the environment and other agents adapt. \parencite[p.~94]{page2011diversity} Gerald Gaus argues, similarly, that once an open society reaches a certain level of complexity, it cannot be planned or even fully understood, and one cannot accurately predict which of “the indefinite possibilities” will be realized. \parencite[p.~14]{gaus2021open}

\section{Large justificatory spaces}

A central claim of this thesis is that public justification aims to connect reasons to coercive law in a way that is publicly intelligible and stable enough to underwrite legitimacy over time. Complexity introduces a distinct obstacle: it expands the \emph{space} of feasible institutional instantiations so dramatically that it becomes difficult for citizens to know what, exactly, their reasons commit them to---and to know whether a given policy remains justified as circumstances shift.

Page offers a simple way to see why complex systems generate massive possibility spaces: if a category of objects has $D$ dimensions and each dimension has $X$ options, then the “size of the possible” is $X^{D}$. \parencite[p.~130]{page2011diversity} Political institutions and policy regimes are (i) highly dimensional (many margins of design, implementation, enforcement, and interpretation) and (ii) high in cardinality (many feasible choices per margin). Even if public reason fixes high-level constitutional values, the remaining design space is enormous.

Moreover, the structure of complex systems is not merely large; it is interactive. Holland stresses that complex adaptive systems are characterized by “intricate hierarchical arrangements of boundaries and signals,” and that, despite abundant description, “we still know little about how to steer these systems.” \parencite[p.~1]{holland2012signals} Steering is hard in part because outcomes depend on conditional interactions: “the behavior of the aggregate is not simply the sum of the agent actions.” \parencite[p.~2]{holland2012signals} In political terms, this means that even well-justified reforms can have emergent consequences that are difficult to anticipate from the justificatory inputs alone.

To connect this to justification: a public reason argument often proceeds as if one can (a) identify a normatively salient target (say, fair equality of opportunity) and (b) select an institutional mechanism that realizes it under roughly stable empirical assumptions. But as Page’s landscape framework makes vivid, the mapping from design to payoff can be rugged: landscapes may have “lots of peaks,” such that local improvements do not guarantee global improvement. \parencite[p.~98]{page2011diversity} And when landscapes “dance,” the problem is harder still: the peaks move as other parts of the system adapt. \parencite[p.~94]{page2011diversity}

This helps clarify why complexity is not reducible to the \emph{fact of reasonable pluralism}. Pluralism increases the diversity of reasons and evaluative perspectives; complexity increases the number of \emph{institutional realizations} compatible with any fixed set of high-level reasons. That is why complexity increases the possibility of a gap between (i) endorsing a political value in abstracto and (ii) knowing what concrete institutional commitments follow from that endorsement.

\section{Exogenous drift: shocks, technology, and changing environments}

The first source of drift is exogenous: changes in background conditions that are not primarily produced by the policy under evaluation (e.g., technological change, demographic shifts, pandemics, geopolitical crises). In the landscape idiom, exogenous drift changes the payoff function.

Page defines \emph{dancing landscapes} as coupled rugged landscapes that “fluctuate due to adaptations by other species,” capturing co-evolutionary processes. \parencite[p.~94]{page2011diversity} The political analogue is that the consequences of institutional choices fluctuate because the social environment to which those institutions apply changes. Exogenous drift is not merely an empirical inconvenience; it alters what counts as an effective realization of political values, thereby altering the justificatory status of policies that were once reasonable.

Gaus’s analysis of the open society strengthens this point. He argues that the open society is characterized by “autocatalytic diversity,” an ever-increasing range of diversity propelling yet more diversity, and that once a society reaches a certain level of complexity it cannot be fully understood or planned. \parencite[p.~14]{gaus2021open} Later he emphasizes that the functioning and consequences of social rules are “increasingly opaque,” and that diversity and complexity are themselves autocatalytic, generating “runaway complexity.” \parencite[p.~102]{gaus2021open} If this is right, then exogenous drift is not episodic; it is a standing feature of open societies.

In such an environment, justificatory arguments that rely on relatively stable empirical generalizations are vulnerable. Even if the underlying political values remain fixed, the evidential basis for policy appraisal shifts as the social world changes. Complexity thus turns the stability problem into something like a moving-target problem: the reasons do not simply “apply” to a stable set of options; the option space and its causal structure evolve.

\section{Endogenous drift: polarization dynamics and institutional feedback}

The second source of drift is endogenous: the system changes in response to itself. Political institutions shape incentives, identities, trust relations, and information flows; these, in turn, reshape the political environment to which justification must respond.

Vallier’s diagnosis of contemporary democracies is a helpful case study of endogenous drift. He argues that Americans are finding it harder to trust one another, that social and political trust have fallen dramatically, and that partisan distrust now extends to ordinary voters, not merely political elites. \parencite[p.~1]{vallier2020trust} He also emphasizes how partisan identity increasingly drives other social identities, such that “politics is swallowing everything else,” and he links falling trust to increased polarization. \parencite[p.~3]{vallier2020trust} The philosophical point is not merely that polarization is bad; it is that polarization dynamics can \emph{reconfigure} the justificatory environment from within, altering what citizens can reasonably accept as public reasons and altering how they interpret the same political values.

Ostrom’s account of polycentric governance makes the same structural point in a different idiom. She argues that, given the complexity of real-world field settings, we should abandon “one size fits all” policies and instead fit institutional rules to specific social-ecological settings; configurational approaches are needed to study what enhances or detracts from “the emergence and robustness of self-organized efforts within multilevel, polycentric systems.” \parencite[p.~642]{ostrom2010polycentric} This is an endogenous-drift claim insofar as institutional arrangements are themselves part of the environment: as rules change, they generate new strategic patterns, new equilibria, and new conflicts, which then become inputs to further rule change.

Complex adaptive systems are, in Holland’s terms, composed of networks of adaptive agents whose aggregate behavior influences agent behavior “as well as vice-versa.” \parencite[p.~2]{holland2012signals} Endogenous drift is therefore unsurprising: laws and institutions are not applied to passive subjects; they are applied to adaptive agents who learn, coordinate, and resist. When justificatory arguments treat institutional effects as approximately stable, they overlook that governance is itself part of the causal web.

\section{Non-linear drift: thresholds, tipping points, and phase transitions}

Complexity also makes drift \emph{non-linear}. Even when change is gradual at the micro level, macro-level consequences can be discontinuous: tipping points, cascades, and phase transitions.

Page’s threshold models provide a vivid illustration. He shows how variation in thresholds, when combined with positive feedback, can create tipping points such that a system that looks stable can suddenly “tip.” \parencite[p.~162]{page2011diversity} On his telling, this is one reason why predicting the timing and speed of riots is notoriously difficult, and why major political transformations can occur rapidly once they begin. \parencite[pp.~162--163]{page2011diversity}

He generalizes the point via the language of phase transitions. As systems approach a phase transition, the relevant landscape can change shape: what was a peak can become a flat spot, selective pressure weakens, and variation increases prior to major change. \parencite[p.~163]{page2011diversity} The political analogue is that apparent stability (e.g., stable partisan coalitions, stable support for constitutional norms) may conceal fragility: small shifts in information, trust, or institutional constraints can precipitate rapid transformation.

Holland’s emphasis on conditional interactions reinforces this non-linearity: because aggregate behavior is not simply additive, small changes in signals or boundaries can have large and sometimes surprising downstream effects. \parencite[p.~2]{holland2012signals} Once we move from linear causal models to interactive networks, it becomes much harder to infer long-run stability from short-run equilibrium.

The upshot for public reason is sobering. If justificatory stability depends on (i) a relatively predictable mapping from institutional choice to political outcomes and (ii) a relatively stable public culture that can interpret those outcomes, then phase-transition dynamics threaten both. The landscape can shift abruptly, and so can the social meaning of policies.

\section{What complexity does to Rawlsian stability}

Rawls’s political liberalism seeks not merely a ceasefire among competing doctrines, but “stability for the right reasons”: a constitutional regime that citizens can endorse from within their reasonable comprehensive doctrines, not merely as a modus vivendi. \parencite[p.~459]{rawls2005political} Public reason, on this view, is not simply a rhetorical restraint; it is part of a moral account of legitimacy.

Rawls emphasizes that the content of public reason is given by a family of liberal political conceptions of justice that (i) apply to the basic structure, (ii) can be presented independently of comprehensive doctrines, and (iii) can be worked out from fundamental ideas implicit in a constitutional regime’s public political culture (citizens as free and equal, society as fair cooperation). \parencite[p.~453]{rawls2005political} In other words, Rawls’s stability story depends on a public political culture that can supply shared justificatory materials.

Complexity disrupts this picture in at least three ways.

First, complexity enlarges the policy realization space so dramatically that citizens can disagree (reasonably, and even within a shared political conception) about which realizations actually instantiate their shared values. Page’s analysis of high-dimensional possibility spaces underscores why. \parencite[p.~130]{page2011diversity}

Second, complexity makes the consequences of social rules increasingly difficult to understand, predict, and evaluate. Gaus’s claim that in the open society the functioning and consequences of social rules become “increasingly opaque,” with diversity and complexity generating runaway complexity, helps explain why citizens may fail to see how the political values they endorse relate to the policies they confront. \parencite[p.~102]{gaus2021open}

Third, complexity turns political justification into a moving-target problem: the landscape dances. Page’s claim that dancing landscapes reflect fluctuating payoffs due to co-evolution, and that on dancing landscapes “the peaks change,” implying the need for continued exploration even after one has found the better option, provides a model for how justificatory equilibrium is undermined by continual change. \parencite[pp.~94, 123]{page2011diversity}

Rawls is not blind to fragility. He explicitly notes that harmony and concord among doctrines and a people’s affirming public reason are “not a permanent condition of social life,” but depend on “the vitality of the public political culture” and on citizens’ being devoted to realizing the ideal of public reason; citizens can “easily fall into bitterness and resentment” and come to ignore it. \parencite[p.~485]{rawls2005political} Complexity provides a structural explanation for why such vitality is hard to sustain: the more the social world becomes an opaque, high-dimensional, co-evolving system, the harder it is for citizens to experience political justification as (i) intelligible, (ii) stable, and (iii) responsive to their reasons over time.

\section{Negative feedback loops and partial stabilizers}

Complexity does not only produce drift; it can also generate stabilizing mechanisms. The question for public reason is whether such mechanisms can be harnessed to preserve legitimacy without collapsing into technocratic control or abandoning the aspiration to stability for the right reasons.

Holland’s methodological lesson is instructive. He contrasts naive trend-based prediction with the identification of mechanisms---Bjerknes’ “fronts” in meteorology---that tell us where to look for relevant data and how to make conditional predictions. \parencite[p.~2]{holland2012signals} The analogy is that, even in complex political systems, institutional design can sometimes identify and strengthen stabilizing mechanisms (monitoring, contestation, error-correction) rather than attempting to fully predict or centrally plan outcomes.

Page gives a simple example of stabilizing negative feedback. In a threshold model of beach attendance (and in an analogous example of bees regulating hive temperature), uniform thresholds can create volatile oscillation, while variation in thresholds dampens swings and promotes stability. \parencite[p.~161]{page2011diversity} This suggests a political analogue: diversity across jurisdictions or institutions (different thresholds, different modes of response) can sometimes stabilize a system by preventing synchronized overreactions and by distributing exploratory effort across the space of possible reforms.

Ostrom’s defense of polycentric governance extends this insight into political economy. Because complex settings require institutional fit, and because self-organization can be more robust under multilevel, polycentric arrangements, a system of multiple overlapping centers of decision can operate as a discovery process and an error-correction process. \parencite[p.~642]{ostrom2010polycentric} Polycentricity can also constrain drift by localizing mistakes: when governance is not fully centralized, failures can be detected and corrected without imposing uniform error on the entire system.

Finally, Gaus argues that an open society depends on a shared, rule-based morality: “we cannot have” a dynamic open society “without it.” \parencite[p.~2]{gaus2016rulebased} In a heterogeneous open society, shared rules can function as a stabilizing boundary condition: they enable coordination among agents who disagree about the good, and they make behavior more predictable even when reasons diverge. \parencite[p.~3]{gaus2016rulebased} Such rule-based constraints are not a complete substitute for public justification, but they can reduce the volatility of justificatory drift by anchoring expectations.

\section{Taming drift: toward a complexity-aware Rawlsianism}

If drift is a standing condition of open societies, then the stability problem should be reframed. Instead of seeking a once-and-for-all equilibrium of public reasons and institutional forms, a complexity-aware view should seek institutions and norms that (i) keep drift within acceptable bounds, (ii) preserve the intelligibility of public justification, and (iii) maintain citizens’ allegiance to political values even as the environment changes.

Rawls’s own distinctions help here. Stability for the right reasons requires more than using public reason as a peacekeeping device; it requires “firm allegiance” to political ideals and values. \parencite[p.~459]{rawls2005political} And Rawls stresses that the durability of public reason depends on the vitality of public political culture and citizens’ devotion to the duty of civility. \parencite[p.~485]{rawls2005political} These claims can be interpreted as institutional and cultural design constraints: a polity must cultivate and maintain the social bases of public justification, not assume them as background facts.

Gaus’s account of the open society suggests a complementary shift in aspiration: political philosophy should treat complexity and deep diversity as a \emph{parameter}, not as a variable that can be engineered away. \parencite[p.~15]{gaus2021open} If justification takes place within the adjacent possible of existing institutions, then legitimacy should be understood as an achievement of an adaptive order, not as the implementation of an ideal blueprint.

Three design commitments follow, each consistent with the sources reviewed here:

\begin{enumerate}[leftmargin=*, itemsep=0.5em]
  \item \textbf{Institutional diversification and experimentation.} Because the policy space is vast and the landscape dances, a polity should preserve room for exploration and revision rather than assuming that exploitation (locking in a single institutional solution) is always rational. \parencite[p.~123]{page2011diversity}
  \item \textbf{Polycentric feedback and local fit.} Because one-size-fits-all rules fail in complex settings, and because self-organization can be robust under multilevel polycentric arrangements, legitimacy-enhancing reforms should often be structured to encourage local adaptation, monitoring, and learning. \parencite[p.~642]{ostrom2010polycentric}
  \item \textbf{Rule-based baselines and shared expectations.} Because heterogeneity undermines reliance on shared virtue or thick ethical consensus, open societies require rule-based moral and legal constraints that stabilize expectations and coordinate behavior, thereby helping sustain the conditions in which public justification can function. \parencite[pp.~2--3]{gaus2016rulebased}
\end{enumerate}

These commitments do not eliminate drift; they aim to make drift governable. They also underscore why legitimacy cannot be secured only by supplying correct reasons at a single time. In complex, autocatalytically diverse societies, maintaining legitimacy requires sustaining the institutional ecology that makes public justification intelligible and stable enough over time.

\section{Summary}

Complexity adds a structural limit to public reason. It expands the space of feasible institutional instantiations, makes steering difficult due to conditional interactions, and introduces exogenous, endogenous, and non-linear drift into the justificatory landscape. Page’s dancing landscapes model the moving-target character of policy evaluation; Holland explains why steering is hard; Gaus explains why open societies become increasingly opaque and unpredictable; Ostrom shows why institutional fit and polycentric feedback matter; Vallier illustrates how trust and polarization dynamics can reshape the justificatory environment from within; and Rawls provides the normative target of stability for the right reasons as well as an acknowledgment of the fragility of public political culture. \parencite[pp.~14, 102]{gaus2021open} \parencite[pp.~94, 123, 162--163]{page2011diversity} \parencite[pp.~1--2]{holland2012signals} \parencite[pp.~1, 3]{vallier2020trust} \parencite[pp.~642--643]{ostrom2010polycentric} \parencite[pp.~453, 459, 485]{rawls2005political}

\noindent The next chapter can build on this analysis by asking what institutional and epistemic supports are needed for public reason to remain action-guiding under conditions of drift, and how those supports interact with the thesis’s earlier arguments about cognition and mediated politics.

\chapter{Attempt at Recovery and a Calibrated Rebuttal}

The preceding chapters argued that Rawlsian legitimacy-through-public-justification is not merely a normative aspiration; it is an aspiration with \emph{feasibility conditions}. If stability ``for the right reasons'' is meant to do justificatory work, then citizens must be able to \emph{track} whether coercive outcomes remain publicly justifiable across time, and the justificatory ``landscape'' must be stable enough for that tracking to succeed. This chapter pauses before drawing the strongest skeptical conclusion. It reconstructs the most charitable, best-case reading of Rawlsian political liberalism and then develops the strongest internal recovery package available in the attached literature.

The structure is deliberately opponent-friendly. I first present a Rawlsian picture on which stability is not a one-shot miracle of convergence but a dynamic property maintained by institutions and public political culture. That picture already suggests a division of justificatory labor that can reduce individual cognitive demands. I then turn to the most sophisticated refinement of this strategy in the attached sources: Weithman’s argument that Rawls’s late ``pluralism about justice'' concession pushes stability away from a \emph{single} fixed point and toward a \emph{set-valued} notion of equilibrium. \parencite[pp.~205--206]{Weithman2023FixedPoints} Finally, I reply using the timescale and granularity framework developed earlier: even after the strongest coarse-graining move is granted, tracking and landscape stability problems do not disappear. They relocate.

\section{The best-case Rawlsian picture}

Rawls frames public reason as a feature of a \emph{well-ordered constitutional democracy}: the idea of public reason ``belongs to a conception of a well-ordered constitutional democratic society.'' \parencite[p.~765]{rawls1997public} The point of the ideal is not that citizens suddenly share a comprehensive moral doctrine. Instead, Rawls starts from a permanent feature of free institutions: \emph{reasonable pluralism}. Under the ``culture of free institutions,'' citizens come to affirm irreconcilable comprehensive doctrines, and they ``cannot reach agreement or even approach mutual understanding'' on that basis. \parencite[p.~766]{rawls1997public} The role of public reason is to specify how citizens can nonetheless relate to one another as free and equal when constitutional essentials and matters of basic justice are at stake. \parencite[pp.~767, 772]{rawls1997public}

A Rawls-friendly reading can therefore resist a crude interpretation of stability as a demand for high-resolution, citizen-level, synchronic convergence. The stability Rawls wants is a diachronic property of a regime’s basic structure and public political culture: it is the capacity of liberal institutions to reproduce themselves in a way that citizens can endorse for the right reasons. Weithman emphasizes this dynamic orientation when he notes that Rawls’s treatment of stability promises help in understanding how liberal institutions can \emph{reproduce themselves} under non-ideal conditions like ours. \parencite[p.~1187]{Weithman2024Rescuing} If that is the aim, then stability should not be read as a stationary point that politics must forever occupy; it should be read as an ongoing pattern of self-maintenance in the face of changing conditions.

It is also natural---and not unfaithful to Rawls’s setting---to connect this reading to general lessons about complex adaptive systems. Holland notes that complex adaptive systems (including governments and markets) exhibit ``intricate hierarchical arrangements of boundaries and signals'' and that, despite abundant description, ``we still know little about how to steer these systems.'' \parencite[p.~1]{holland2012signals} Taken as an analogy rather than an attribution, the Rawls-friendly thought is: a well-ordered liberal society is not a simple machine that runs once set in motion. It is a complex system whose stability depends on institutions that channel contestation, transmit signals of justificatory failure, and correct course without requiring each citizen to model the full political environment at fine resolution.

This sets up the first steelman of the ``tracking'' critique: perhaps Rawls does \emph{not} require citizens to track everything. If public reason is primarily about constitutional essentials and matters of basic justice, then what citizens must track may be coarse-grained---roughly, whether the regime remains within the bounds of a liberal-democratic constitutional order. \parencite[pp.~767--768]{rawls1997public} On this reading, much of the justificatory work is done by institutions and practices that citizens can monitor at a higher level of abstraction.

\section{Recovery route 1: division of justificatory labor}

Rawls explicitly builds an institutional division of justificatory labor into the idea of public reason. The public political forum includes ``the discourse of judges'' (especially in a supreme court), ``the discourse of government officials,'' and ``the discourse of candidates for public office.'' \parencite[p.~767]{rawls1997public} Ordinary citizens are not absent from the picture, but their role is partly mediated: public reason is institutionalized through constitutional interpretation, legislative debate, and electoral accountability. \parencite[pp.~767--768]{rawls1997public} Under a charitable interpretation, Rawls is not asking each citizen to personally compute whether each law is justified at each time; he is describing a \emph{public practice} of justification in which different institutions shoulder different epistemic tasks.

This route becomes more vivid when contrasted with Habermas’s schematic models of democracy. Habermas characterizes the liberal view as one on which the political process is structured by strategic competition for power; voters’ decisions resemble market choices that express preferences and license access to office. \parencite[p.~3]{habermas1994three} Even within that liberal frame, the constitution is supposed to \emph{channel} competition through basic rights and the separation of powers. \parencite[p.~7]{habermas1994three} Rawls’s institutional emphasis can be read as an attempt to specify \emph{how} such channeling should proceed when legitimacy is tied to justification: public reason is meant to regulate what counts as an admissible political argument when constitutional essentials are at stake. \parencite[pp.~767--768]{rawls1997public}

Habermas’s discourse-theoretic model pushes the same division-of-labor idea further. On his view, democratic procedures and their communicative presuppositions serve as ``sluices'' through which public use of reason can rationalize administrative decisions; the public sphere operates as a ``far-flung network of sensors'' that registers social problems and generates influential opinions. \parencite[p.~9]{habermas1994three} What results is not the abolition of administration but a reorientation: communicative power cannot ``rule'' by itself, yet it can ``point'' administrative power in specific directions. \parencite[p.~9]{habermas1994three} Read alongside Rawls, the thought is not that institutions replace citizens, but that the epistemic burdens of political justification are \emph{distributed} through publicly structured forums that citizens can contest and revise.

This division-of-labor strategy is, at minimum, a serious answer to a simplistic version of the tracking objection. If courts, legislatures, and public reasoning norms perform some of the tracking work, then citizens can maintain democratic control without individually maintaining a high-resolution map from reasons and facts to policy consequences at every moment. The ideal can aim for a publicly shared justificatory \emph{practice}, not a heroic cognitive achievement by isolated individuals.

However, the strategy has a latent tension. The more justificatory labor is offloaded onto institutions and experts, the more citizens must rely on \emph{trust} in those institutions; yet Rawlsian legitimacy is supposed to be grounded in citizens’ ability to see coercive law as justified in terms they can share. The very mechanism that makes the ideal feasible threatens to weaken the citizen-facing uptake that gives the ideal its distinctive moral point. This tension does not defeat the strategy, but it motivates a second Rawls-friendly route: the appeal to well-orderedness and the educative effects of institutions.

\section{Recovery route 2: well-ordered liberal society and educative institutions}

Rawls’s stability story is not only institutional but also psychological and cultural: a well-ordered society is one in which citizens are motivated to comply with just institutions and to support them over time. Weithman’s reconstruction makes clear that Rawls’s \emph{argument} for stability relies on the educative effects of institutions. In Weithman’s terms, Rawls’s earlier equilibrium claim (for justice as fairness) depended on the idea that institutions satisfying justice as fairness shape citizens’ values in ways that support their continued reproduction. \parencite[p.~206]{Weithman2023FixedPoints} The point, again, is dynamic: institutions are not merely constrained by citizens’ prior values; they also cultivate dispositions and expectations that can stabilize cooperation.

Weithman’s more recent defense of Rawls against Cohen underscores why this matters. Cohen criticized Rawls for connecting stability and justice, arguing that stability is ``alien'' to justice and should be studied separately. \parencite[p.~1187]{Weithman2024Rescuing} But Weithman argues that Rawls’s treatment of stability can be defended, precisely because it illuminates how liberal institutions reproduce themselves under non-ideal conditions. \parencite[p.~1187]{Weithman2024Rescuing} Even if one ultimately agrees with Cohen that justice and stability are distinct questions, Rawls’s project plainly treats stability as part of what a political conception must address if it is to guide a constitutional democracy marked by pluralism.

This educative route also connects Rawls’s project to a broader deliberative-democratic ambition. Bohman describes the deliberative-democratic tradition as emerging from critique of standard liberal-democratic practice and as aiming to vindicate democratic legitimacy through public reasoning rather than mere aggregation of preferences. \parencite[pp.~400--401]{bohman1998comingage} Yet, as Bohman emphasizes, deliberative democracy ``comes of age'' partly by grappling with feasibility: it must articulate institutional forms and social conditions under which public reasoning can play the legitimating role democratic theory assigns to it. \parencite[pp.~400--403]{bohman1998comingage} A Rawls-friendly strategy can therefore claim that the well-ordered society is not an ad hoc assumption; it is the name for a package of institutions, norms, and civic education that makes citizen uptake of public reason feasible.

If that package were robust, the tracking burden might be reduced in two ways. First, citizens could treat certain public institutions---courts, legislatures, a free press, civic organizations---as reliable \emph{signals} about whether the regime remains within liberal bounds. Second, the ongoing operation of those institutions could stabilize the interpretive environment in which citizens form political judgments, thereby slowing the drift that undermines justificatory beliefs. The possibility of such ``taming'' is not implausible. The question is whether it is sufficient once we grant the strongest refinement of the Rawlsian picture: the shift from point-valued to set-valued stability.

\section{Weithman’s recovery route 3: from fixed points to set-valued stability}

The most sophisticated version of the coarse-graining reply in the attached sources is Weithman’s argument that Rawls’s later writings force a structural revision in how we think about stability.

\subsection{Rawls’s concession and the need for a new equilibrium concept}

Weithman argues that Rawls made an ``important concession'' about the likely target of an overlapping consensus: it is more likely to focus on ``a class of liberal conceptions'' than on a single conception such as justice as fairness. \parencite[p.~205]{Weithman2023FixedPoints} Let that class be $\{L_{1},\ldots,L_{n}\}$, where each $L_{i}$ is a Rawlsian liberal conception (satisfying the familiar conditions about basic liberties, reciprocity, and the like). \parencite[p.~205]{Weithman2023FixedPoints} Weithman’s key interpretive move is that Rawls’s concession is not a small sociological aside; it changes what stability can reasonably mean. If a society can remain well-ordered while being effectively regulated by different members of $\{L_{1},\ldots,L_{n}\}$ over time, then stability cannot be identified with occupying one particular point forever. \parencite[pp.~205--206]{Weithman2023FixedPoints}

\subsection{Stability as ``remaining within a ball''}

To make the point vivid, Weithman uses the metaphor of a ``ball'' in a space of conceptions or institutional realizations. He suggests interpreting Rawls’s remarks about a ``narrow'' range and a ``center'' of a focal class in geometric terms: the acceptable liberal conceptions may be thought of as points within a ball centered on justice as fairness. \parencite[p.~205]{Weithman2023FixedPoints} The crucial revision is that stability can be understood as remaining \emph{within} that ball over time, even if the regime moves around inside it. \parencite[p.~205]{Weithman2023FixedPoints} This is the most sophisticated version of the granularity reply: citizens need not track exact convergence to a single conception; they need to track that the regime stays within an acceptable neighborhood.

\subsection{From point-valued to set-valued equilibrium}

Weithman argues that if Rawls’s concession is taken seriously, Rawls must endorse a new equilibrium claim: societies effectively regulated by Rawlsian liberal conceptions are in equilibrium. \parencite[p.~206]{Weithman2023FixedPoints} But that equilibrium concept cannot be point-valued. If a society can transition from a point identified as just by one member of $\{L_{1},\ldots,L_{n}\}$ to a point identified as just by another member (or by a hybrid) while remaining well-ordered, then equilibrium must be \emph{set-valued}. \parencite[p.~206]{Weithman2023FixedPoints} In Weithman’s terms, Rawls’s concession therefore requires not only a new idea of well-orderedness but also ``a new concept of equilibrium'' and a new argument for its existence. \parencite[p.~206]{Weithman2023FixedPoints}

One payoff of this reconstruction is that it clarifies why ``dynamic'' critics of Rawls sometimes think Rawlsian theory is inevitably static. Weithman summarizes the dynamic theorist as focusing on a justice \emph{process} in which values are endogenous and agents’ considered judgments reflexively adapt to the structures they help to produce; dynamic theory therefore ``eschew[s] the assumption of fixed point equilibria of justice.'' \parencite[p.~207]{Weithman2023FixedPoints} Weithman’s response is not to deny that institutions shape values; Rawlsians agree on that. Rather, the dispute turns on whether the values encouraged by just institutions can be sufficiently stable to ``equilibrate'' those institutions, or whether institutional rules must continually evolve to accommodate evolving values. \parencite[p.~207]{Weithman2023FixedPoints} The set-valued equilibrium picture is a Rawls-friendly way to concede movement without abandoning the hope that movement is bounded in a way that can support legitimacy.

\subsection{Why this matters for the thesis}

For present purposes, Weithman’s reconstruction is the strongest internal defense against a feasibility critique that targets citizen tracking. If stability is set-valued, then what citizens must track can be coarser. They need not identify a unique, fine-grained conception as \emph{the} publicly justified one. Instead, they can reasonably hope that political life remains within a range of liberal conceptions that are each acceptable as bases of cooperation among free and equal citizens.

There is also a natural connection to the general phenomenon that diversity is both productive and risky in complex systems. Page emphasizes that diversity can provide insurance, improve productivity, spur innovation, enhance robustness, and produce collective knowledge, while also warning that diversity ``is no panacea'' and can contribute to breakdown. \parencite[p.~3]{page2011diversity} The Rawlsian hope, in Weithman’s form, is that allowing a class of liberal conceptions (rather than a single doctrine-like point) preserves some benefits of pluralism while still bounding political contestation within liberal limits. This is exactly what a feasibility-sensitive Rawlsian should want.

The question, then, is not whether this move makes Rawls’s project \emph{more} plausible. It does. The question is whether it makes the project plausible \emph{enough} once we explicitly foreground the timescale and granularity parameters.

\section{A calibrated rebuttal: what still fails after the strongest coarse-graining move}

\subsection{What Weithman’s move genuinely fixes}

Weithman’s reconstruction blocks a straw objection. If a critic assumes that Rawlsian stability requires convergence on a single doctrine-like conception of justice that must be endorsed in essentially the same way by all reasonable citizens, the critic targets an implausibly demanding picture. Rawls’s later writings, as Weithman reads them, explicitly allow a plurality of liberal conceptions; stability need not mean fixity at one point. \parencite[pp.~205--206]{Weithman2023FixedPoints} This is a genuine improvement in the feasibility profile of Rawlsian legitimacy: it relaxes what must be shared (granularity) and allows variation through time (a dynamic rather than static image).

\subsection{What still fails: tracking and landscape stability relocate rather than disappear}

Even if stability is set-valued, citizen tracking does not evaporate. It changes target. Instead of tracking whether society is regulated by \emph{this} particular conception, citizens must track whether society remains \emph{within} the liberal class $\{L_{1},\ldots,L_{n}\}$ (or within the ``ball'') \emph{and} whether the institutions that mediate justification are reliably keeping political outcomes within those bounds. \parencite[p.~205]{Weithman2023FixedPoints} That is still a substantive epistemic task. In Rawls’s own institutional picture, the public political forum includes officials and judges whose discourse is supposed to be governed by public reason. \parencite[p.~767]{rawls1997public} But citizens must still be able to judge (in an appropriately public way) when those institutions are functioning as intended and when they are drifting, captured, or strategically re-describing illiberal policies as members of the liberal family.

The same relocation occurs for landscape stability. The set-valued strategy relies on there being a stable-enough mapping from (i) endorsement of a liberal class of conceptions and (ii) institutional mediation of those conceptions into (iii) concrete policy outcomes that are recognizably within liberal bounds. Yet Holland stresses that complex adaptive systems are characterized by intricate signal-and-boundary hierarchies and that we know little about steering them. \parencite[p.~1]{holland2012signals} In such systems, it is difficult to ensure that coarse-grained commitments (``stay within the liberal ball'') reliably generate fine-grained outcomes that remain inside the intended boundary, especially when incentives and strategic behavior encourage boundary-pushing.

Habermas’s discourse-theoretic picture provides a helpful way to articulate this vulnerability. If the public sphere is a network of sensors that reacts to problems and ``stimulate[s] influential opinions,'' and if public opinion can only \emph{point} administrative power in specific directions, then the translation from communicative power into administrative outcomes is mediated and fallible. \parencite[p.~9]{habermas1994three} That fallibility is compatible with democracy; it is part of what makes politics politics. But it means that the Rawlsian must specify a regime of institutions and public norms under which (a) the sensors reliably detect justificatory failures and (b) the administrative system reliably responds in a way that preserves liberal boundaries \emph{over the relevant timescale}.

This is where the timescale bite becomes unavoidable. If the timescale $\tau$ over which stability-for-the-right-reasons is demanded is short, then citizens must update their beliefs about whether they remain within the liberal ball at a high frequency; the epistemic environment must supply rapid, reliable signals of boundary crossings. If $\tau$ is long, then the cognitive demands may be reduced, but at a cost: legitimacy judgments become increasingly retrospective, and the ideal risks becoming a story we tell after the fact about institutions that happened to remain acceptable, rather than a condition citizens can reasonably rely on while authorizing coercion.

The upshot is not that the Weithman-style coarse-graining move is empty. It is the most compelling answer to the simplest form of the feasibility critique. But it does not remove the need for citizen-facing tracking; it shifts it to a higher-order problem: tracking \emph{membership conditions} for the liberal family (what counts as ``within the ball'') under interpretive drift and strategic manipulation, and tracking whether institutional mediation is preserving the boundary between liberal and illiberal outcomes.

\subsection{Calibrated conclusion: a conditional feasibility thesis}

A fair conclusion must be conditional rather than dismissive. Weithman’s defense of Rawls against Cohen is instructive here: Weithman argues that Rawls’s arguments for the stability of a just society are more limited and tentative than Rawls acknowledged, and that locating those limits shows ``the points at which unjust societies such as our own need to be sh[ored up].'' \parencite[p.~1187]{Weithman2024Rescuing} That is exactly the right posture for the present thesis. The strongest Rawlsian recovery package does not trivialize feasibility concerns; it clarifies the conditions under which the ideal might be realized.

On the best-case picture developed in this chapter, Rawlsian stability-through-public-justification is feasible only under a restricted set of dynamic and epistemic conditions: a political culture that sustains public reason, institutions that distribute justificatory labor while remaining publicly contestable, and a sufficiently stable mapping from liberal commitments to policy outcomes over the relevant timescale. If those conditions fail, Rawlsian legitimacy does not merely become harder to achieve; it loses grip as a practical ideal of citizen authorization, because the public reasons that are supposed to ground legitimacy cannot be reliably tracked as the political environment shifts.

\section{Objections and Replies}

Two objections deserve explicit acknowledgement, because they sharpen the transition to the constructive institutional discussion that follows.

First, a Rawlsian might insist that public reason is an \emph{ideal} of civic respect, not a sociological prediction, and therefore that feasibility concerns are beside the point. But Rawls’s own framing locates public reason within a well-ordered constitutional democracy and assigns it the job of specifying how citizens can be bound to honor constitutional essentials under reasonable pluralism. \parencite[pp.~765--767]{rawls1997public} If stability-for-the-right-reasons is meant to underwrite legitimacy in that setting, then feasibility constraints are internal to the ideal’s role: a practice that cannot be sustained by the agents to whom it applies cannot secure the legitimacy it is invoked to secure.

Second, a deliberative democrat might object that too much institutional mediation risks turning democratic self-government into administration constrained by a thin constitutional morality. Habermas’s contrast between liberal, republican, and discourse-theoretic models highlights this worry: the liberal model treats politics as strategic competition, whereas the discourse-theoretic model gives the public sphere a constitutive role in rationalizing administration. \parencite[pp.~3, 9]{habermas1994three} Bohman likewise emphasizes that deliberative democracy’s maturation involves specifying institutions that make public reasoning effective without collapsing into elitism or mere preference aggregation. \parencite[pp.~400--403]{bohman1998comingage} The present thesis accepts the spirit of this objection: democratic legitimacy cannot be replaced by opaque expertise. But the feasibility problem remains even for deliberative democrats. If public reasoning is to play the legitimating role assigned to it, then citizens must have public, contestable means of tracking justificatory status over time.

This is the pivot to the next chapter. If the strongest Rawlsian recovery package still leaves a residual tracking-and-drift vulnerability, then the natural response is not to abandon democratic justification. It is to ask what kinds of institutions can function as publicly inspectable epistemic supports---institutions that reduce the cognitive burden on individuals while preserving the democratic roles of criticism, revision, and authorization.

\chapter{Institutions as Cognitive Prosthetics}

\section{Motivation: from individual tracking to institutional scaffolding}

The central worry running through this thesis is not that citizens \emph{lack} reasons, but that the justificatory demands of political liberalism can be undermined when citizens cannot \emph{track} the reasons that are purported to justify coercive laws---especially under conditions of scale, complexity, and persistent disagreement.
Rawls’s later statement of political liberalism is explicit that the idea of public reason \enquote{belongs to a conception of a well ordered constitutional democratic society} and is \enquote{part of the idea of democracy itself} precisely because modern democratic citizenship is marked by the \enquote{fact of reasonable pluralism}. \parencite[p.~765]{rawls1997public}
If that is right, then public justification is not an optional moral adornment; it is a constitutive feature of legitimate democratic authority.

But Rawls also insists that public reason has a structure that includes not only officials’ justificatory burdens but also citizens’ role in \emph{checking} whether coercive norms satisfy the criterion of reciprocity. \parencite[p.~767]{rawls1997public}
In a representative democracy, citizens are to \enquote{think of themselves as if they were legislators} and ask what statutes, supported by reciprocity-satisfying reasons, they would most reasonably enact. \parencite[p.~769]{rawls1997public}
This idealized picture assigns ordinary citizens a distinctive epistemic responsibility: they must be able to assess, at least in broad outline, whether the public reasons offered for foundational political decisions are sincere, relevant, and not merely rhetorical.

The difficulty is that the modern public sphere does not reliably supply the informational and deliberative conditions needed to make this role feasible. 
Even when we bracket disagreement about final values, the justificatory premises for large-scale policy are often technical, probabilistic, and mediated through complex institutions.
Without supportive infrastructure, the \emph{tracking problem} becomes a collective-action problem: each person has limited time and attention, while the benefits of careful scrutiny are diffuse and the incentives for distortion can be concentrated.
If political liberalism is committed to public justification, then it must also be committed to the practical possibility of making public justification \emph{publicly available}.

Jonathan Quong captures the point in a crisp slogan: \enquote{public justification requires public reasons}, that is, reasons that can be recognized as valid political considerations by those to whom justification is owed. \parencite[p.~10]{quong2011liberalism}
The present chapter argues that this requirement cannot be met by exhortation alone.
What is needed is an institutional strategy: shifting some of the epistemic and justificatory burden from individual citizens to public, contestable, and revisable forms of \emph{institutional scaffolding}.
In Weithman’s recent defense of Rawlsian stability, the aim is framed in explicitly political terms: to understand how liberal democratic institutions might \emph{reproduce themselves} even \enquote{under non-ideal conditions such as our own}. \parencite[p.~1188]{Weithman2024Rescuing}
Institutional design, on this view, is not a technocratic add-on; it is part of what it means to take seriously the internal demands of a democratic conception of legitimacy.

This chapter develops one way to think about that design task. 
The central proposal is that certain institutions can function as \emph{cognitive prosthetics}: they do not replace citizens’ judgment, but they make it easier for citizens to exercise that judgment responsibly by improving the epistemic environment in which public reasons are produced, tested, and maintained.
The goal is to make the practice of public reason more robust to scale, complexity, and strategic behavior, without collapsing into technocracy or manipulation.

\section{Institutions as cognitive prosthetics: definition and scope}

Call an institution a \emph{cognitive prosthetic} when it systematically reduces citizens’ epistemic burdens in ways that (i) preserve citizens’ status as free and equal co-authors of law, and (ii) make publicly justifiable reasons easier to \emph{identify}, \emph{evaluate}, and \emph{revise} over time.
The metaphor is intentionally modest.
A prosthetic does not guarantee success; it expands feasible agency under constraints.
Likewise, institutional prosthetics do not guarantee truth, consensus, or justice.
They aim to make the public justificatory economy less brittle by embedding procedures that generate reliable public signals and correction mechanisms.

A helpful starting point comes from Hart’s account of law as \enquote{the union of primary and secondary rules}. \parencite[pp.~79--80]{hart2012concept}
Hart argues that a social order of primary rules alone is vulnerable to three systematic defects: \emph{uncertainty} about what the rules are, the \emph{static} character of rules that cannot be deliberately adapted, and the \emph{inefficiency} of diffuse social pressure as a mode of enforcement and dispute settlement. \parencite[pp.~92--93]{hart2012concept}
The legal solution is not (only) better moral character; it is a different \emph{institutional layer}. 
Secondary rules provide public criteria for identifying valid rules (\enquote{rules of recognition}), procedures for deliberate legal change (\enquote{rules of change}), and authoritative mechanisms for resolving disputes (\enquote{rules of adjudication}). \parencite[pp.~94--97]{hart2012concept}

Hart’s framework generalizes beyond jurisprudence.
A political order that demands public justification is vulnerable to parallel defects.
Without public criteria for what counts as an acceptable public reason, justification becomes uncertain.
Without public procedures for revising justificatory settlements, public reason becomes static, even as circumstances change.
Without forums that can authoritatively register challenge and correction, justificatory practices become inefficient or collapse into informal power.
Seen in this light, public reason is not only a set of norms for discourse; it is a candidate for institutionalization through rules and procedures that structure how reasons are offered, assessed, and updated.

This institutional perspective is also central to Buchanan and Tullock’s account of constitutional choice.
They emphasize that judgments about collective action depend on the decision rules governing it, and that meaningful evaluation requires attention to the \enquote{constitutional} level at which those rules are chosen. \parencite[p.~6]{buchanan1962calculus}
Because agreement on ultimate rules is costly, individuals may nonetheless find it advantageous to \enquote{agree in advance to certain rules} that structure subsequent collective choice. \parencite[p.~7]{buchanan1962calculus}
Whatever one thinks of their broader project, the key point for present purposes is this: citizens can exercise agency \emph{over time} by choosing and revising the institutional procedures that later shape their political options.
Cognitive prosthetics belong primarily at this constitutional and quasi-constitutional level---the level at which citizens can reasonably demand that justification be made visible, contestable, and revisable.

\section{Institutional filters and feedback loops}

If cognitive prosthetics are to support public reason, they must operate less like censorship and more like \emph{filters} and \emph{feedback loops}.
A filter, in this sense, is a procedure that improves the signal-to-noise ratio of public justification by making reasons easier to inspect and compare.
A feedback loop is a mechanism that enables public justification to learn from error, abuse, or changed circumstances.
This is not a call to replace politics with expertise; it is a call to build institutions that help democratic politics remain publicly justifiable under conditions that otherwise predict rhetorical drift.

Two bodies of work provide complementary guidance about what such institutions should look like.

\subsection{Polycentric structure as an epistemic check}

Ostrom’s account of polycentric governance begins from a diagnosis of complex social systems: \enquote{one size fits all} policies are not effective because institutional rules must be adapted to specific social-ecological contexts. \parencite[p.~642]{ostrom2010polycentric}
The concept of the polycentric is then defined in explicitly institutional terms: it \enquote{connotes many centers of decision making that are formally independent of each other}, where jurisdictions may nonetheless function coherently through patterned interaction. \parencite[p.~643]{ostrom2010polycentric}
Polycentric arrangements, on this view, are not merely decentralized for its own sake. They enable experimentation, learning, and mutual monitoring across sites of decision.

For public reason, the relevance is direct.
When justificatory claims are complex and the costs of error are high, a single centralized epistemic authority becomes a single point of failure.
Polycentricity offers a distinctive form of epistemic redundancy: multiple bodies can evaluate the same justificatory questions from different vantage points, compare outcomes, and force public reasons to survive cross-checking rather than rhetorical momentum.
Ostrom’s empirical discussion of policing illustrates how large centralized systems can fail to outperform smaller units embedded in local contexts, and how multiple units can generate performance gains through comparability and responsiveness. \parencite[pp.~644--645]{ostrom2010polycentric}
For our purposes, the point is structural: institutions that support public reason should be multiply sited, overlapping, and capable of mutual scrutiny.

\subsection{Legibility, thin simplifications, and the need for local knowledge}

James Scott’s diagnosis of modern statecraft is, in one sense, the negative image of Ostrom’s constructive argument.
Scott describes the premodern state as largely \enquote{blind} to local society and explains modern state-building as an effort to make society \enquote{legible} through administrative ordering and standardization. \parencite[pp.~2--3]{scott1998seeing}
These state simplifications can be indispensable for public goods and rights; but they can also be dangerously \enquote{thin} representations of social reality.
When simplifications are treated as if they were the whole truth, they become tools for coercive redesign rather than instruments of limited governance.

Scott’s clearest warning is that the most tragic episodes of state-initiated social engineering have tended to combine four elements: state simplifications; a high-modernist ideology of rational design; the coercive power of an authoritarian state; and a weakened civil society unable to resist. \parencite[pp.~4--5]{scott1998seeing}
From the standpoint of public reason, this is a cautionary tale about institutional prosthetics.
Institutions that aim to improve public justification can, if badly designed, become instruments for manufacturing consent, narrowing perceived options, and enforcing schematic ideals in the name of rational administration.

The antidote is not to abandon institutional design, but to build it around \emph{feedback} from the practical knowledge that resists simplification.
Scott emphasizes that large social processes are typically more complex than any schema, and that failures of planned \enquote{scientific} ordering are often mitigated (or quietly corrected) by informal improvisations outside the plan. \parencite[p.~309]{scott1998seeing}
If public reason is to travel from theory to practice, its supporting institutions must therefore be structured to detect when justificatory schemas have become detached from lived social realities, and to allow correction before coercive error becomes entrenched.

\subsection{Collective intelligence without technocracy}

If Scott warns against schematic administration and Ostrom argues for polycentric experimentation, Landemore offers a positive epistemic case for democratic inclusion.
Her core thesis is that, under conditions conducive to deliberation and majority rule, democracy can be a better decision procedure than non-democratic alternatives such as councils of experts or benevolent dictators. \parencite[p.~3]{landemore2013democratic}
The key mechanism is that including more people in decision-making tends to increase \emph{cognitive diversity}, a crucial ingredient of collective intelligence in both problem solving and prediction. \parencite[pp.~3--4]{landemore2013democratic}

This matters for cognitive prosthetics because it shifts the design target.
The point is not merely to inform citizens \emph{about} expert judgments; it is to structure institutions so that expertise is \emph{integrated into} democratic contestation and widened perspectives.
Institutions that support public reason should thus avoid substituting a single epistemic hierarchy for public judgment.
Instead, they should create public pathways by which diverse perspectives can test, reinterpret, and (when appropriate) reject elite framings.

\subsection{What institutional filters can look like}

At a high level of abstraction, institutional filters and feedback loops that support public reason will tend to include the following components:

\begin{enumerate}
    \item \textbf{Public standards of justificatory relevance.} Analogous to Hart’s rule of recognition, the system needs visible criteria for what counts as a public political reason in the relevant domain. \parencite[pp.~94--95]{hart2012concept}
    \item \textbf{Procedures for revision.} Justificatory settlements must be revisable under publicly specified rules of change rather than only through crisis or informal power. \parencite[pp.~95--96]{hart2012concept}
    \item \textbf{Authoritative venues for challenge.} Disputes about whether justificatory standards have been met require accessible procedures and institutions with the standing to register reasons as defeated, insufficient, or inconclusive. \parencite[p.~97]{hart2012concept}
    \item \textbf{Polycentric evaluation.} Multiple independent bodies should be empowered to assess justificatory claims, compare conclusions, and force public reasons to withstand cross-checking. \parencite[p.~643]{ostrom2010polycentric}
    \item \textbf{Local feedback against thin simplifications.} Institutional processes should be designed to surface local knowledge and practical experience that schematic accounts miss, thereby preventing the ossification of legibility-based governance. \parencite[pp.~2--3, 309]{scott1998seeing}
\end{enumerate}

These are not blueprints; they are design constraints.
They state what a public-reason-supporting institutional ecology must be able to do if it is to reduce epistemic burdens without replacing democratic agency.

\section{Paternalism and manipulation: the hard objection}

The most serious moral objection to cognitive prosthetics is that they may smuggle in a form of paternalism: instead of persuading citizens, institutions might manage them.
This worry has two aspects.

First, institutional scaffolding might implicitly treat citizens as epistemically incompetent, thereby undermining their standing as free and equal.
Second, even if institutions aim at epistemic improvement, they might be captured by elites and used to steer the public toward preferred outcomes through agenda control, selective disclosure, or rhetorical framing.

Quong’s discussion of paternalism helps clarify what is at stake.
He suggests that one reason the harm principle is attractive is that it helps ensure citizens are not treated \enquote{in a certain type of paternalistic fashion, one which demeans their moral status as free and equal persons}. \parencite[p.~56]{quong2011liberalism}
Paternalistic coercion can express domination or disrespect by signaling that some citizens cannot be trusted to decide for themselves. \parencite[p.~58]{quong2011liberalism}
If the institutional prosthetics proposed here expressed such contempt, they would be incompatible with political liberalism rather than a remedy for it.

Rawls’s account of public reason points to a different direction.
Public reason is the \enquote{reason of free and equal citizens}, and its nature and content must be \enquote{public}, expressed in public reasoning that aims to satisfy reciprocity. \parencite[p.~767]{rawls1997public}
Officials are to \enquote{explain to other citizens their reasons} for supporting fundamental positions, and citizens are to hold officials to this duty through a \enquote{duty of civility}. \parencite[p.~769]{rawls1997public}
This relationship is incompatible with institutional strategies that rely on secrecy, manipulation, or unreviewable expertise.
If institutions are to support public reason, they must amplify the visibility and contestability of reasons, not bypass them.

Scott’s warnings make the capture problem vivid.
State simplifications can become instruments of coercive redesign when allied to high-modernist ideology and concentrated state power. \parencite[pp.~4--5]{scott1998seeing}
In that setting, what begins as an attempt to rationalize governance can end as a project of authoritarian domination.
The paternalism objection is therefore not a peripheral worry; it is an ever-present failure mode for any institutional strategy that aims to improve the epistemic environment.

\section{Design principles: democratic safeguards for epistemic institutions}

To avoid paternalism and manipulation, cognitive prosthetics must be designed as \emph{democratic safeguards} rather than epistemic command centers.
The basic idea is to make epistemic assistance compatible with democratic agency by ensuring that the institutions remain transparent, contestable, and polycentric---and by locating them within a constitutional framework that citizens can reasonably endorse.

\subsection{Transparency as a condition of public justification}

Because public justification requires public reasons, justificatory institutions must produce outputs that are legible to citizens and available for scrutiny. \parencite[p.~10]{quong2011liberalism}
Rawls’s insistence that public reason is \enquote{public} in its very nature means that justificatory premises cannot be treated as proprietary expertise.
Institutional reports, standards, and recommendations must be accompanied by explanations that connect technical claims to shared political values and make uncertainty explicit. \parencite[p.~767]{rawls1997public}

\subsection{Contestability through structured challenge and adjudication}

Transparency without contestability is fragile: it allows powerful actors to \emph{display} reasons without being responsive to criticism.
Hart’s account of adjudication captures why institutionalized challenge matters. 
Rules of adjudication confer powers to make authoritative determinations about whether rules have been breached and define procedures for doing so. \parencite[p.~97]{hart2012concept}
Analogously, a public-reason-supporting institutional ecology needs standing venues in which justificatory failures can be named and registered: where citizens and civil society actors can contest whether public reasons meet reciprocity and whether policies rest on thin simplifications or strategic distortion.

\subsection{Polycentricity and anti-capture redundancy}

Ostrom’s definition of polycentricity highlights a further safeguard: multiple formally independent centers of assessment reduce the risk that any single body becomes an epistemic monopolist. \parencite[p.~643]{ostrom2010polycentric}
Polycentric design provides redundancy against capture and encourages learning through comparison.
In the justificatory domain, this implies overlapping review bodies with different appointment procedures, different constituencies, and different methodological defaults---so that public reasons must travel through multiple filters rather than a single gatekeeper.

\subsection{Local knowledge and humility against thin simplifications}

Scott’s account of legibility suggests another principle: justificatory institutions must be humble about the limits of schematic knowledge and must systematically incorporate feedback from those affected by policy. \parencite[pp.~2--3]{scott1998seeing}
Because large-scale schemes often fail when they treat simplified representations as if they were the whole truth, institutions should treat contestation and local testimony not as noise but as epistemic resources. \parencite[p.~309]{scott1998seeing}
This is especially important for policies that reorganize lived social practices (e.g., housing, education, policing), where practical knowledge is often dispersed and context-sensitive.

\subsection{Constitutional authorization and citizen control over time}

Finally, the paternalism objection is mitigated when epistemic institutions are grounded in constitutional procedures that citizens can endorse and revise.
Buchanan and Tullock stress that individuals can rationally agree in advance to decision rules that structure later collective choice, even when agreement is costly and substantive outcomes uncertain. \parencite[pp.~6--7]{buchanan1962calculus}
For cognitive prosthetics, the analogous thought is that citizens can reasonably authorize institutions that improve the justificatory environment \emph{provided} those institutions remain accountable and revisable.
Citizen agency is preserved not only in day-to-day participation but also in the ongoing power to reform the rules that govern justification.

\section{Bridge: stability, justice, and Weithman’s rescue project}

Weithman’s defense of Rawls against Cohen begins from the idea that stability is not merely a sociological add-on but is intimately connected to justice as Rawls understands it. \parencite[p.~1188]{Weithman2024Rescuing}
Yet Weithman also emphasizes the difficulty of applying Rawls’s stability analysis under non-ideal conditions and the importance of understanding how liberal democratic institutions might reproduce themselves in societies like our own. \parencite[p.~1188]{Weithman2024Rescuing}

The cognitive-prosthetics framework proposed here is one way to take that challenge seriously.
If stability partly depends on citizens’ dispositions to affirm and sustain just institutions, then the epistemic environment in which citizens evaluate political power matters.
Institutions that make public reasons more visible, contestable, and resilient can support the kind of civic confidence needed for stability \emph{without} demanding implausible levels of individual competence.
At the same time, because these institutions are themselves sites of power, they must be constrained by the very values they aim to serve: reciprocity, publicity, and respect for citizens as free and equal.

The result is demanding but not utopian.
It does not assume ideal citizens; it demands institutions that help actual citizens do what democratic citizenship already asks of them: to assess, contest, and (when warranted) endorse the exercise of coercive political power in terms that others can reasonably accept. \parencite[p.~769]{rawls1997public}
In this sense, institutional prosthetics are not a departure from public reason but an attempt to build the conditions under which public reason can plausibly function as a regulative ideal in real political life.

\section{Summary and transition}

This chapter argued that the tracking problem motivating the thesis should be understood as a structural challenge to public justification under modern conditions.
The response is not to lower the justificatory standard, nor to replace citizens with experts, but to develop an institutional ecology that functions as cognitive prosthetics: public standards, polycentric review, structured contestation, and feedback from local knowledge.
Grounded in the values of public reason, such institutions can make legitimacy more robust while respecting the anti-paternalist core of political liberalism.

In the concluding chapter, I draw together the thesis’s main claims and clarify what it would mean, politically and institutionally, to pursue this project in societies marked by polarization, strategic misinformation, and persistent pluralism.

\chapter{Conclusion: What Follows from the Invisible Limits}

\section{What this thesis has shown}

The central claim of this thesis has been that legitimacy-through-public-reason depends on two feasibility conditions that are easier to invoke than to satisfy. Citizens must be able to track whether coercive outcomes remain justifiable over time, and the political environment must be stable enough for that tracking to be more than episodic. The conclusion of the argument is not that public justification is morally unimportant. It is that the standard Rawlsian route from public reason to long-run stability is much more environmentally and epistemically demanding than its proponents often acknowledge.

The sources considered in the later chapters help clarify why. Gaus argues that the open society is not simply a liberal order marked by diversity; it is 
characterized by \enquote{autocatalytic diversity}, an ever-expanding range of differences that generates yet more differences, such that, beyond a certain level of complexity, the system \enquote{cannot be planned or even fully understood}. \parencite[pp.~14--15]{gaus2021open}
That observation matters for a thesis about Rawlsian stability because it means that the justificatory environment is not fixed. Political agents are not reasoning once and for all in a settled world. They are reasoning in a world whose categories, alignments, and stakes continue to shift.

Gaus also shows, from within public reason liberalism itself, that the problem is not exhausted by simple disagreement over first principles. Even when an abstract freestanding justification is available, it may still be too indeterminate to yield the fine-grained social morality required for ordinary collective life, since once people bring their \enquote{full range of values, religious convictions, and other views} to bear, they may continue to affirm the abstract conclusion while disagreeing sharply about its interpretation. \parencite[p.~42]{gaus2011order}
And when public justification is modeled more realistically, members of the public evaluate proposals not by appeal to one shared stock of reasons but on the basis of \enquote{mutually intelligible, but diverse, evaluative standards}. \parencite[p.~283]{gaus2011order}
The upshot is that indeterminacy is not an accidental deviation from public justification. It is one of its standing conditions in a diverse order.

The empirical studies attached to this chapter underscore how quickly this diverse environment is reshaped by media and social exposure. Mutz argues that exposure to dissimilar political views can expand citizens' awareness of the rationales for opposing positions and can increase political tolerance, especially when cross-cutting exposure links people to views held by those with whom they have meaningful associations. \parencite[pp.~111--114]{mutz2002crosscutting}
Yet Levendusky and Malhotra show that media coverage portraying the electorate as deeply polarized increases perceived polarization, moderates respondents' own issue positions, and heightens negative affect toward the opposing party. \parencite[pp.~283--284, 291--293]{levendusky2016media}
The same public sphere can therefore generate two contrary tendencies at once: exposure to disagreement can sometimes enlarge tolerance, but mediated representations of polarization can also intensify antagonism. A theory of legitimacy that depends on durable public uptake cannot ignore that instability in the conditions of uptake.

Taken together, these materials support the thesis's main conclusion. The difficulty with Rawlsian stability is not merely that citizens sometimes disagree, or that institutions sometimes malfunction. It is that the work of maintaining public justification must be done in a political world marked by indeterminacy, adaptive change, strategic mediation, and shifting perceptions of division. Under those conditions, the burden of diachronic tracking is far more severe than standard idealized accounts suggest.

\section{What this thesis has \emph{not} shown}

That conclusion should be stated carefully. It does not follow that ideal theory is worthless, that normative political philosophy must collapse into political sociology, or that demanding principles are somehow philosophically discredited by human failure. Estlund's \emph{Utopophobia} is especially useful here. He insists that \enquote{an account that requires more than we ever expect to achieve is not thereby in the least flawed}, even if questions about whether we should set out for it are separate and practically important. \parencite[p.~4]{estlund2019utopophobia}
Likewise, he argues that being unrealistic may be a defect of a \emph{proposal}, but not of a \emph{principle}; \enquote{being unrealistic is not a defect in a principle of justice, though it would be in a proposal}. \parencite[p.~12]{estlund2019utopophobia}
And he repeatedly distinguishes the philosophical task of identifying standards from the practical task of deciding whether, how, and under what conditions those standards should guide institutional action. \parencite[pp.~4, 10--12]{estlund2019utopophobia}

Accordingly, this thesis has not shown that Rawls's ideals are false simply because they are difficult to realize. Nor has it shown that political philosophy must never outrun present feasibility. Its argument is narrower and more internal. When a theory makes legitimacy depend on a specifically public and stable relation between reasons, institutions, and citizen endorsement, feasibility constraints enter the theory from the inside. They are not alien \enquote{realist} objections imposed from outside the normative project; they are part of what must be true if that legitimating route is to work.

For the same reason, the thesis has not shown that democratic contestation is futile, or that citizens are incapable of learning. The empirical literature attached to this chapter points in the opposite direction: citizens can become more aware of opposing rationales under cross-cutting exposure, and they can react in nuanced ways to mediated polarization. \parencite[pp.~111--114]{mutz2002crosscutting}; \parencite[pp.~291--293]{levendusky2016media}
The problem is not that democratic agents are inert. It is that the environments in which they reason are unstable enough, and the tasks assigned to them demanding enough, that we should hesitate before building legitimacy around the expectation of durable, high-fidelity justificatory uptake.

\section{Implications for Rawlsians}

For Rawlsians, the most immediate implication is not abandonment but restriction and clarification. If public justification is to continue doing legitimating work, Rawlsians need a more explicit account of the domain in which stability for the right reasons can plausibly be expected, and a more candid account of the institutional supports required to sustain it.

The attached sources suggest why. Gaus's criticism of high-level abstraction is directly relevant: even robust freestanding justification may fail to provide the specificity needed for actual social-moral rules, once full evaluative diversity is admitted into the justificatory process. \parencite[p.~42]{gaus2011order}
Moreover, the publicity ideal itself is more demanding than simple formal acceptability. Gaus argues that strong publicity requires a situation in which all normal agents have sufficient accessible reasons to endorse the rules and can, in a free moral equilibrium, regard others as similarly situated. \parencite[p.~297]{gaus2011order}
That is an ambitious condition. It cannot be assumed merely because a reasonable political conception has been articulated at the level of theory.

The conclusion to draw is not that Rawlsians must retreat to a purely philosophical picture detached from practice. It is rather that they must decide among three possibilities. First, they may narrow the scope of the legitimacy claim, treating stability for the right reasons as a regulative aspiration appropriate only under relatively favorable conditions. Second, they may accept a looser, more mediated notion of citizen uptake, according to which public justification depends less on each citizen's detailed tracking than on the reliability of institutions that render reasons publicly available and corrigible. Third, they may concede that the full Rawlsian ambition is best understood as an evaluative standard rather than a generally feasible account of political legitimacy.

Nothing in this trilemma refutes justice as fairness as a moral ideal. But it does indicate that late-Rawlsian arguments about stability cannot be insulated from the political epistemology of justification. Once the surrounding environment becomes more visibly dynamic, polarized, and mediated, the justificatory work that stability is supposed to do requires more institutional support than the classical statement of the view tends to foreground.

\section{Implications for convergence theorists}

The post-Rawlsian convergence tradition offers one prominent strategy for lowering the demands of consensus-based public reason. Vallier's dissertation formulates this strategy in especially clear terms. He argues that an alternative conception of public reason liberalism should combine a \enquote{convergence conception of public reasons} with a \enquote{moderate conception of idealization}. \parencite[p.~8]{vallier2011diss}
On this view, reasons offered to justify coercion \enquote{need not be shared}; convergence therefore relaxes the standard requirement that all justified coercion rest on the same publicly accessible reasons. \parencite[p.~8]{vallier2011diss}

This move undeniably accomplishes something important. Vallier argues that convergence is \enquote{far more permissive} than consensus because it \enquote{recognizes all intelligible reasons as justificatory}, and for that reason it avoids many of the integrity costs that burden mainstream public reason liberalism. \parencite[p.~191]{vallier2011diss}
In that respect, convergence offers a genuine answer to one familiar criticism of Rawlsian restraint.

But convergence does not make the tracking problem disappear. It relocates it. Once justification can proceed through diverse but mutually intelligible reasons, the task of political evaluation shifts from identifying one shared justificatory basis to navigating a wider field of acceptable reasons, defeaters, and idealizing assumptions. Vallier himself acknowledges this when he turns to idealization. A viable moderate idealization, he says, must avoid both \enquote{gross errors in reasoning and information} and the excessive normalization of citizens' reasons; it must attribute reasons in ways consistent with agents' deeply held values while also specifying standards of relevance, collection cost, and processing cost. \parencite[pp.~255--257]{vallier2011diss}
That is a sophisticated and important concession. It recognizes that any public justification view still needs a theory of what citizens may reasonably be taken to know, infer, and process.

Gaus's account of public justification helps make the point even sharper. If members of the public evaluate proposals through diverse but mutually intelligible standards, and if the move from abstract principles to concrete rules is pervaded by indeterminacy, then convergence does not simplify democratic navigation so much as pluralize it. \parencite[pp.~42, 283]{gaus2011order}
The number of justificatory pathways expands, but the problem of locating oneself within them remains. In that respect, convergence theorists alleviate one burden of public reason while preserving, and in some ways intensifying, another: the burden of determining when a coercive outcome remains publicly justified amid multiple intelligible routes to endorsement.

The implication is not that convergence fails. It is that convergence theory, no less than Rawlsian theory, needs a more explicit political epistemology. If the justificatory space is widened, then the institutions that help citizens interpret, compare, and contest justificatory claims become more important, not less.

\section{Implications for democratic design}

For democratic design, the attached sources point toward a different lesson: pluralism should be processed, not wished away. Institutions should not be evaluated by whether they produce a final consensus on all contested questions, but by whether they allow citizens to encounter disagreement, register reasons, and revise political judgments without turning ordinary conflict into permanent enmity.

The empirical literature provides a useful starting point. Mutz's work suggests that cross-cutting exposure can have genuine democratic benefits by increasing awareness of opposing rationales and fostering political tolerance. \parencite[pp.~111--114]{mutz2002crosscutting}
This finding matters because it shows that institutions and networks that expose citizens to disagreement are not merely costs to be managed; they can be conditions of democratic learning. But Levendusky and Malhotra show that exposure is not enough in the abstract. How disagreement is framed matters profoundly. When media coverage depicts the electorate as polarized and divided, citizens come to perceive greater mass polarization, they move their own issue positions toward the center, and yet they simultaneously become more negatively disposed toward the opposing party. \parencite[pp.~283--284, 291--293]{levendusky2016media}
Design, then, cannot be indifferent to mediation. Institutions must attend not only to whether citizens encounter disagreement, but to the social forms in which disagreement is rendered salient.

Gaus's later work pushes the design implication beyond media effects to the level of institutional architecture. Because the open society is marked by autocatalytic diversity and by complexity that cannot be fully planned or mastered, normative political philosophy must work within what he calls the \enquote{adjacent possible} or the \enquote{neighborhood} of existing institutions rather than imagining that the whole social order can be deliberately redesigned from above. \parencite[pp.~14--15]{gaus2021open}
More concretely, he argues that effective governance is most plausible where agents share a \enquote{pressing problem-solving context}: a common perception of a problem, rough agreement about a range of plausible solutions, and a belief that resolving the problem is better than leaving it unresolved. Under those conditions, heterogeneity is not abolished, but it is reduced enough to make coordinated problem-solving feasible. \parencite[pp.~199--200]{gaus2021open}

This line of argument helps explain why the constructive turn in this thesis has emphasized institutional scaffolding rather than institutional omniscience. In a highly diverse order, democratic institutions should help citizens solve specific problems under conditions of contestability, rather than promise total reconciliation at the level of a complete public philosophy. Gaus's warning about multiple policy goals strengthens the point. As governors take on more and more simultaneous targets, the complexity of the system reasserts itself because the output of one policy becomes an input into others. \parencite[p.~218]{gaus2021open}
A political order that demands too much global coherence from its citizens will therefore tend to reproduce the very overload it cannot handle.

There is also a cautionary implication. Gaus is skeptical of frameworks that try to secure harmony by splitting society into thick, like-minded communities under a thin liberal shell, since fragmentation of that kind can reduce the desire for reconciliation and increase conflict. \parencite[pp.~206--207]{gaus2021open}
That point matters for contemporary democratic design. The answer to pluralism is not purified enclave democracy. It is the design of institutions that can stabilize cooperation across disagreement without pretending that disagreement will disappear.

\section{Future directions}

If the invisible limits identified in this thesis are real, then three lines of future work become especially important.

First, the empirical study of democratic environments should focus more directly on \emph{drift regimes}: the rate, scale, and form by which mediated political contexts alter what citizens take to be publicly relevant, publicly reasonable, and politically threatening. Mutz and Levendusky point in this direction already. Together they show that the political consequences of disagreement depend on the channels through which disagreement is encountered. \parencite[pp.~111--114]{mutz2002crosscutting}; \parencite[pp.~291--293]{levendusky2016media}
A fuller account of legitimacy under pluralism will require more precise maps of those channels.

Second, work in public reason liberalism should continue to integrate diversity and indeterminacy rather than treat them as residual complications. Gaus's analyses of abstract justification, intelligible but diverse standards, and strong publicity suggest that public justification has always contained a more complex structure than the shared-reasons picture admits. \parencite[pp.~42, 283, 297]{gaus2011order}
Vallier's convergence liberalism and moderate idealization show one route forward, but they also reveal the need for more explicit standards concerning citizens' informational burdens and institutional supports. \parencite[pp.~8, 255--257]{vallier2011diss}

Third, the debate about ideal theory should move away from the sterile opposition between moral aspiration and practical realism. Estlund's distinction between principles and proposals is the right starting point. \parencite[pp.~10--12]{estlund2019utopophobia}
The question is not whether political philosophy may think beyond current feasibility. It is whether a given theory presents its normative aspirations as action-guiding proposals, as evaluative standards, or as a legitimacy-conferring account whose own success depends on facts about institutions and citizens. Greater clarity on that point would improve both the philosophical and the institutional sides of the debate.

\section{Final reflections}

The argument of this dissertation began with a question about public reason, but it ends with a broader claim about the task of political philosophy in complex democracies. The task is not simply to identify the most elegant justificatory ideal and wait for the world to catch up. Nor is it to surrender normativity to the brute facts of polarization, misinformation, and conflict. It is to ask what must be true for a justificatory ideal to function as a practice of political life, and what institutional forms can help free and equal citizens navigate a world in which disagreement is durable, environments drift, and social orders are more complex than any single model can contain.

That task remains philosophical. But it is philosophical in a way that must remain answerable to political environments as they are encountered by actual agents. If the invisible limits traced here are real, then the future of public justification depends less on discovering a final consensus than on designing forms of democratic life that can sustain mutual accountability under conditions of enduring plurality. That is not the end of political philosophy. It is a demanding account of where its next work lies.

\appendix

\chapter{A minimal formal sketch}
\label{app:formal-model}

This appendix provides a deliberately minimal formalization of the toy model in Section~\ref{sec:toy-model}. It is not meant to prove a theorem about Rawlsian legitimacy. It is meant to clarify what, exactly, the \emph{tracking} requirement asks citizens to do.

\section{States, policies, and justificatory status}

Let $S_t$ denote a \emph{state} of social and institutional facts at time $t$ (economic structure, administrative practice, interpretive doctrine, salient empirical regularities). Let $P$ denote a set of policy options. Let an outcome function $O$ map pairs $(p,S_t)$ to outcome profiles $O(p,S_t)$.

Let $J$ be a (public-reason) justificatory rule that assigns a status to policies in states:
\[
J(p,S_t) \in \{0,1\},
\]
where $J(p,S_t)=1$ means \enquote{publicly justifiable (reciprocity-satisfying) in $S_t$} and $J(p,S_t)=0$ means \enquote{not publicly justifiable in $S_t$}. This is a crude stand-in for Rawls's richer idea that justification must be acceptable to citizens under a criterion of reciprocity. \parencite[pp.~446--447]{rawls2005political}

\section{Citizen models and bounded updating}

Citizen $i$ does not observe $S_t$ directly. Instead, the citizen receives noisy signals $x_t$ generated by the environment. The citizen forms an internal model $\widehat{S}^i_t$ and an estimate of justificatory status $\widehat{J}^i_t(p)$.

A predictive-processing interpretation suggests two features:

\begin{enumerate}[leftmargin=*, itemsep=0.4em]
    \item \textbf{Update via error:} citizens revise $\widehat{S}^i_t$ when signals generate sufficient prediction error. \parencite[p.~816]{friston2005theory}
    \item \textbf{Compression:} citizens represent $S_t$ in a lower-dimensional or sparse code (a summary statistic) rather than as a full description. \parencite[pp.~491--492]{ganguli2012compressed}
\end{enumerate}

Write the citizen's compression as a map $g_i$ from states to a low-dimensional representation $z^i_t$:
\[
z^i_t = g_i(S_t), \quad \text{with } \dim(z^i_t) \ll \dim(S_t).
\]
The citizen's justificatory estimate is therefore a function of the compressed representation:
\[
\widehat{J}^i_t(p) = \widehat{J}^i(p, z^i_t).
\]

\section{Drift and tracking error}

Model drift as change in the underlying state:
\[
S_{t+1} \neq S_t,
\]
with the important case being \emph{small but persistent} drift: $\|S_{t+1}-S_t\|$ is small at each step but accumulates over time.

Define citizen $i$'s tracking error for policy $p$ at time $t$ as:
\[
\epsilon^i_t(p) = \left| J(p,S_t) - \widehat{J}^i_t(p) \right|.
\]

The qualitative claim defended in the chapter can be stated as:

\begin{quote}
If (i) drift changes the justificatory status of some policies by altering $O(p,S_t)$ and its public meaning, and (ii) citizens update $\widehat{S}^i_t$ mainly when error signals are salient, while (iii) their representations are compressed, then $\epsilon^i_t(p)$ can remain small for a time and then grow abruptly once accumulated drift crosses the citizen's error-detection threshold.
\end{quote}

This is a schematic way of expressing the chapter's core point: JTC presupposes a diachronic alignment relation between citizens' justificatory judgments and the evolving facts that determine whether coercion remains publicly justifiable. Predictive-processing models suggest that, for finite agents, that alignment is fragile under drift.

% ====================
% Bibliography
% ====================

\printbibliography

\end{document}